\chapter{Building Adaptive Hybrid Representations}\label{chap:method}
% ---- TAKEAWAYS ----
% - Joining whole representations adds their irrelevant features along with the useful ones. The hybrid combines and filters at once.
% - A hybrid is an input, not a model: up to a budget of features, selected from several base representations by a search for one dataset and one downstream model.
% - No setting of the search wins on every dataset or model, so importance method, allocation method, representation count, internal evaluator, budget and downstream model are all searched.
% - Every selected feature traces back to its representation, so a blueprint can show a chemist which substructures and descriptors the model relies on.
% - The method needs only a supervised label and several vector representations of the same samples, so it carries over to antibodies, reactions, nucleotides and tumor data.
% - Where one representation already dominates, the hybrid should show little gain. The search is costly and must be rerun when the data, representations or model change.
% ---- END TAKEAWAYS ----

% CODE INPUT
This chapter asks how several representations can be joined without also joining their irrelevant features. We answer with the adaptive hybrid representation, a set of at most a fixed number of features, the budget, selected from several base representations by a search for one dataset and one downstream model. The hybrid is therefore an input to a model rather than a model itself. The search covers the importance method, the allocation method, the number of representations, the internal evaluator, the budget and the downstream model. Because every selected feature traces back to its representation, each finished search yields a blueprint that a chemist can read. The method needs only a supervised label and several vector representations of the same samples, and we test it on small molecules, antibodies, reactions, nucleotide sequences and tumor data. \textbf{[RESULT NEEDED: on how many datasets the hybrid beats the best single base representation, and whether the preferred search settings differ between datasets and models]}.


% TODO: how do I reference these properly? 
\section{Introduction}
% [notes moved to thesis_notes_moved.md]

% CODE INPUT
Chapter~\ref{chap:flaws} showed that standard representations carry many features that do not vary, or do not matter, on a given dataset, and that a distance computed over all of those features inherits the problem. That chapter ended by calling for features chosen for the task from several representations, with the irrelevant ones removed first. Chapter~\ref{chap:robustness} gave a practical reason to search over this choice instead of fixing it, since the best representation changed with both the dataset and the model. The selection described here was first built to compute similarity and was later found to work as a model input.

% CODE INPUT
In this chapter we build the adaptive hybrid representation, hereafter the hybrid, and report first results on when and why it helps. Each hybrid draws on base representations of the same molecules and is specific to one dataset and one downstream model.

% CODE INPUT
This chapter measures prediction only, and distances computed in a hybrid, noise detection and learning with noisy labels are the subject of Chapter~\ref{chap:similarity}. The methods described here are final, and each dataset is split into 5 outer folds. The hybrid beats the best single base representation on \textbf{[RESULT NEEDED: number of the 16 datasets on which the hybrid beats the best single base representation]} of the 16 datasets. \textbf{[RESULT NEEDED: whether the preferred search settings vary between datasets, and which settings]}.

% CODE INPUT
The rest of this introduction reviews ways of combining representations and the alternatives to combining them, then describes the hybrid, its blueprints and its use beyond small molecules, before the methods, results, discussion, limitations and conclusions.


\subsection{Combining representations and the alternatives}
% [notes moved to thesis_notes_moved.md]
% - TODO(lit): read the Turkina2025 full text for whether a fixed number of bits is kept per fingerprint, or all bits are ranked together.

% TODO: review
Prior research in this field shows promise, with conjoint fingerprints achieving stronger performance on binding affinity prediction than their component fingerprints alone across a range of models \citep{Xie2020}. Augmenting structure-based fingerprints like ECFP4 with bioactivity-based fingerprints has also been successful at improving both predictive performance and scaffold-hopping \citep{Laufktter2019}. A very similar fusion strategy, in which features are selected from six non-hashed fingerprints and are combined by feature importance score within individual models, consistently improves predictive accuracy over the original fingerprints \citep{Turkina2025}. 

% TODO: review
The majority of the work in this space deals with combining relatively standard fingerprints and linear representations like SMILES \citep{Xie2020, Laufktter2019}. However, graph-based representations capture important structural information that is often lost in fingerprints. We explore using GNN embeddings within such hybrid fusion methods. We further expand on the mechanisms used to compute feature importance and concatenate features, and we use these hybrids within similarity quantification and noise mitigation tasks. We also explore adding smaller sets of structural features to a single base representation, referred to as augmentation. 

% TODO: REVIEW MODIFICATION
The majority of the work in this space deals with combining relatively standard fingerprints and linear representations like SMILES \citep{Xie2020, Laufktter2019}. However, graph-based representations capture important structural information that is often lost in fingerprints, and we therefore include GNN embeddings in the hybrid. MHG-GNN and GROVER are base representations on every small-molecule property dataset except QM9, and GraphMVP is a base representation on QM9 and LogD (Table~\ref{tab:reps_smallmol}). We further expand on the mechanisms used to compute feature importance and to select and join features. We also use these hybrids for similarity quantification, noise detection and learning with noisy labels. Finally, we explore adding smaller sets of structural features to a single base representation, referred to as augmentation.

% TODO: review and the li and deng defini
% TODO: review, needs manual review
So far, the most common strategy for combining feature sets has been concatenation, which requires only a single downstream model \citep{Stahlschmidt2022}. Concatenation works best when the feature sets are homogeneous and degrades when they are not \citep{Stahlschmidt2022}, and \citet{Deng2023} called the practice ``misleading'' when the goal is to measure the contribution of individual representation types. Late fusion trains a separate model per source and aggregates the predictions \citep{Steyaert2023}, which avoids the dimensionality problem but loses interactions between features from different sources. Recovering those interactions requires intermediate fusion (MOGONET \citep{Wang2021}, cross-attention \citep{Beaude2025}) % TODO: REVIEW MODIFICATION -- was \citep{Zhang2024}, which had no bibliography entry and no known source. CrossAttOmics (Beaude et al. 2025) is a verified cross-attention intermediate-fusion paper. Swap it if you meant a different paper.
, at the cost of end-to-end training on all sources, so that adding a new source means retraining from scratch. Multiple kernel learning \citep{Gonen2011} operates entirely at the kernel level, so no interpretable feature subsets emerge from it. As Steyaert et al.\ noted, there remains no decisive evidence that any of these strategies is broadly superior \citep{Steyaert2023}.

% TODO: REVIEW MODIFICATION
So far, the most common strategy for combining feature sets has been concatenation, which requires only a single downstream model \citep{Stahlschmidt2022}. Concatenation works best when the feature sets are homogeneous and degrades when they are not \citep{Stahlschmidt2022}. It also obscures which representation is responsible for the score. Deng et al.\ found that appending RDKit2D descriptors to GROVER lowered RMSE on 9 of 24 datasets, yet they still called the practice ``misleading'' when the goal is to measure the contribution of individual representation types \citep{Deng2023}. Late fusion trains a separate model per source and aggregates the predictions \citep{Steyaert2023}, which avoids the dimensionality problem but loses interactions between features from different sources. Recovering those interactions requires intermediate fusion (MOGONET \citep{Wang2021}, cross-attention), at the cost of end-to-end training on all sources, so that adding a new source means retraining from scratch. Multiple kernel learning \citep{Gonen2011} operates entirely at the kernel level, so no interpretable feature subsets emerge from it. An example in drug bioactivity prediction is pairwiseMKL \citep{Cichonska2018}. As Steyaert et al.\ noted, there remains no decisive evidence that any of these strategies is broadly superior \citep{Steyaert2023}.

% TODO: review, needs manual review
In practice, most research groups default to one representation and ignore the rest. ECFP is common for small molecules, while deep learning groups may opt for embeddings like MoLFormer or ChemBERTa. ESM-2 embeddings are common for antibodies and can be computed per chain. Some research compares these representations, but little focuses on combining them beyond simple concatenation. Concatenating just three standard sources, however (ECFP4 at 2,048 dimensions, ESM-2 at 1,280, Mordred descriptors at ${\sim}1{,}600$), gives a ${\sim}5{,}000$-dimensional input for a dataset that may contain only a few hundred labeled samples. At these sample-to-feature ratios, more features mean worse predictions \citep{Hughes1968, Bellman1961}. Models sensitive to distance, such as $k$-nearest neighbors (kNN) and SVM, are affected most severely \citep{FernandezDelgado2014}.

% TODO: REVIEW MODIFICATION
In practice, most research groups default to one representation and ignore the rest. ECFP is common for small molecules, while deep learning groups may opt for embeddings like MoLFormer or ChemBERTa. ESM-2 embeddings are common for antibodies and can be computed per chain. Some research compares these representations, but little focuses on combining them beyond simple concatenation. Concatenating just three standard sources, ECFP4 (2,048 features), ESM-2 (1,280) and Mordred descriptors (${\sim}1{,}600$), gives about 5,000 features, whereas a dataset may contain only a few hundred labeled samples. At these sample-to-feature ratios, more features mean worse predictions \citep{Hughes1968, Bellman1961}. Models that compute a distance over every feature, such as $k$-nearest neighbors (kNN) and SVM, are the most exposed, for the reasons given in Chapter~\ref{chap:flaws}. A concatenation of two representations scored between its two parts, rather than above the better one, on \textbf{[RESULT NEEDED: count of datasets]}. We propose that this happens because concatenation adds the irrelevant features of both parts along with the useful ones. A method that joins representations therefore has to combine and filter at the same time.

% TODO: review, needs manual review
The pace of new model development shows no sign of slowing, and extends well beyond small molecules. New models may prefer different feature sets, which this method can provide. Beyond that, new models can be converted into frozen embeddings and used as further input to this framework.

% TODO: REVIEW MODIFICATION
The pace of new model development shows no sign of slowing, and extends well beyond small molecules. A new downstream model may prefer a different feature set, and it therefore receives its own search. A new model can also be converted into frozen embeddings, which enter the method as one more base representation without any change to the method.

% TODO: review, needs manual review
A comparison of 62,820 models in molecular property prediction benchmarks showed that classical RDKit2D descriptors matched or outperformed learned features from graph neural networks and pretrained language models on the majority of datasets \citep{Deng2023}. Descriptor-based models occupied 73\% of top-three positions on 11 benchmarks \citep{Jiang2021}, and in 1,310 ChEMBL assays, feedforward networks on ECFP fingerprints outperformed graph convolutions and SMILES-based recurrent models \citep{Mayr2018}. The input thus often matters more than the model, yet no single representation wins and there is still no systematic way to exploit this. The same pattern holds outside small-molecule chemistry, where no single protein language model consistently achieves the highest performance on all tasks \citep{Kenlay2024, Schmirler2024}. For antibodies, general-purpose models sometimes outperform those built for the domain \citep{Kenlay2024}. In computational pathology, no individual foundation model dominates on clinical endpoints \citep{ElNahhas2025}, and in cancer genomics, integrating additional omics layers can reduce performance when the data types are not properly combined \citep{Steyaert2023}.

% TODO: review, needs manual review
For small molecules alone, the landscape now includes extended-connectivity fingerprints, physicochemical descriptor sets exceeding 1,600 dimensions, graph neural network embeddings \citep{Yang2019, Fang2022}, language models trained on billions of SMILES strings \citep{ross2022}, image-based encodings (molecules rendered as 2D images) \citep{Zeng2022}, and stereoelectronics-infused molecular graphs \citep{Boiko2025}. The same is true for proteins (ESM-2 \citep{Lin2023}, ProtTrans \citep{Elnaggar2022}), antibodies (AntiBERTy \citep{Ruffolo2021}, AbLang \citep{Olsen2022}), histopathology \citep{ElNahhas2025}, and multi-omics cancer profiles \citep{Steyaert2023}. A recent review counted over 40 molecular foundation models published between 2020 and 2024 \citep{Song2026}. Increasingly, new models appear faster than we can compare them \citep{Bommasani2021}, so the choice is not getting simpler. Different sources presumably encode distinct aspects of the same biology. Single data modalities, noisy or incomplete on their own, give more sensitive and robust predictions when combined with independent signals from other modalities \citep{Steyaert2023}. El Nahhas et al.\ found that an ensemble of frozen embeddings from the top four pathology foundation models achieved a mean area under the receiver operating characteristic curve (AUROC) 1.2\% higher than the best individual model on 31 clinical tasks \citep{ElNahhas2025}.

% TODO: REVIEW MODIFICATION
Representations now exist for small molecules, reactions, antibodies, nucleotide sequences and tumor profiles, and those used in this thesis are listed in Tables~\ref{tab:reps_smallmol}, \ref{tab:reps_reactions}, \ref{tab:reps_antibodies}, \ref{tab:reps_nucleotides} and~\ref{tab:reps_tcga}. A recent review counted over 40 molecular foundation models published between 2020 and 2024 \citep{Song2026}. Increasingly, new models appear faster than we can compare them \citep{Bommasani2021}, so the choice is not getting simpler. Different sources encode distinct aspects of the same biology. Single data modalities, noisy or incomplete on their own, give more sensitive and robust predictions when combined with independent signals from other modalities \citep{Steyaert2023}.

% TODO: review, needs manual review
Concatenation, the most common method for combining feature sets, gives the full pool from which the hybrid starts, but on biological sample sizes it can degrade performance instead of helping \citep{Bellman1961, Stahlschmidt2022}. At the kernel level, multiple kernel learning \citep{Gonen2011} (e.g., pairwiseMKL for drug bioactivity \citep{Cichonska2018}) can learn source weights, but no interpretable feature-level selections result. Consensus methods \citep{Valsecchi2020} combine predictions rather than features, and so cannot exploit cases where one source's noise is another's signal. The hybrid instead works at the feature level and selects features that span source boundaries. Perhaps the closest analogy in the optimization literature is the combined algorithm selection and hyperparameter optimization (CASH) problem \citep{Thornton2013} (the joint selection of algorithm and hyperparameters), extended here to include the input encoding as an additional search dimension. Tishby and Zaslavsky \citep{Tishby2015} showed why compression helps, since the information bottleneck principle states that optimal representations maximize mutual information with the target while minimizing mutual information with the input. In effect, the budget forces the model to keep only what matters, and the feature count controls the trade-off between informativeness and dimensionality.

% TODO: REVIEW MODIFICATION
The hybrid search starts from all base representations side by side but works at the feature level. It keeps original features drawn from several representations, and so selects features that span representation boundaries. Perhaps the closest analogy in the optimization literature is the combined algorithm selection and hyperparameter optimization (CASH) problem \citep{Thornton2013}, the joint selection of algorithm and hyperparameters. The hybrid search adds a third choice, namely which features from which representations enter the model.

% TODO: background on feature selection algorithms
% But more of a direct focus on feature selection

% TODO: background on PCA and other dimensionality reduction techniques


% CODE INPUT
Ensembling combines representations after prediction rather than before it, by fitting one model per representation and pooling their predictions. Consensus QSAR follows this approach and combines the predictions of several models instead of their features \citep{Valsecchi2020}, so it cannot exploit cases where one representation's noise is another's signal. In computational pathology, an ensemble built on the four best foundation models reached a mean AUROC 1.2\% above the best single model over 31 clinical tasks \citep{ElNahhas2025}. In an ensemble, however, no single model sees features from two representations together, and every representation and every model must be computed for each new molecule.

% CODE INPUT
Feature selection keeps a subset of the original features and builds no new ones, which separates it from principal component analysis (PCA). Methods are usually grouped into three families \citep{Guyon2003}, which \citet{Saeys2007} reviewed for high-dimensional biological data. Filter methods score features against the labels before any model is fitted, as when minimum-redundancy maximum-relevance selection (mRMR) ranks features by their mutual information with the label minus their redundancy with features already chosen \citep{Peng2005}. Wrapper methods search over subsets of features and score each subset by fitting the model \citep{Kohavi1997}, and in QSAR, descriptors were selected this way with a $k$-nearest-neighbor model \citep{Zheng2000}. Embedded methods select features while the model is fitted, as when the lasso's L1 penalty sets some coefficients to exactly zero \citep{Tibshirani1996}. Recursive feature elimination fits a model, drops the lowest-ranked features and repeats \citep{Guyon2002}. Two pitfalls of feature selection are well known. The selected set can change between resamples of the same data, which the Nogueira index measures across folds \citep{Nogueira2018}, and impurity importance in random forests favors some feature types over others \citep{Strobl2007}.

% CODE INPUT
The hybrid ranks features by the importance scores of a fitted model, as embedded methods do, and then keeps a fixed number of top-ranked features, as a filter does. Classical methods treat all features as one flat list and ignore the representation each feature came from. Their nearest relative to the hybrid is selection over predefined groups, as in the group lasso, which keeps or drops whole groups of variables together \citep{Yuan2006}. The hybrid, in contrast, does not keep or drop whole representations. Four of the five allocation methods in its search give each representation a share of the budget and take that many top-ranked features from it. The fifth, the global pool, gives each representation a share of a candidate pool twice the size of the budget, then ranks the pooled candidates together and keeps the top ones up to the budget. Because only original features are kept, every feature in the hybrid traces back to one representation, which is what makes blueprints possible.

% CODE INPUT
Dimensionality reduction replaces the original features with fewer new features, each a combination of the originals, and in QSAR it turns many correlated descriptors into fewer uncorrelated ones \citep{Niazi2023}. PCA is the standard method for this purpose. Each component is a linear combination of the original features, and components are ordered by the variance they explain \citep{Jolliffe2016}. Because components are usually functions of all the original features, they are hard to interpret \citep{Jolliffe2016}. PCA also depends on the units of the features, which z-scoring each feature first removes \citep{Jolliffe2016}, and this matters when fingerprint bits, descriptors and embeddings are joined. PCA never sees the labels and keeps the directions of largest variance, which need not predict the property, while components with small variance can still predict it \citep{Jolliffe1982}. Partial least squares (PLS) is the supervised counterpart in chemometrics, with components built to covary with the property \citep{Wold2001}. Both methods give mixed components, so a component cannot be traced to one substructure, descriptor or representation, and no blueprint can be read from it. Nonlinear methods include kernel PCA \citep{Niazi2023}, t-distributed stochastic neighbor embedding (t-SNE) \citep{vanderMaaten2008} and uniform manifold approximation and projection (UMAP) \citep{McInnes2018}, although t-SNE and UMAP are used mainly for two- or three-dimensional maps rather than as model inputs.

% CODE INPUT
Every alternative above except ensembling also reduces the number of features. Models may, however, want different numbers of features. Distance-based models such as kNN and SVM use every feature in every distance, whereas tree ensembles choose one feature at each split and can therefore partly ignore irrelevant features. Models may also want different features, and not only different numbers of them. Selection that uses the labels may overfit. Sort \& Slice, which uses no labels, outperformed two label-based schemes for selecting substructure features \citep{sns}, and Dablander suggests that using the labels may add a risk of overfitting \citep[p.~138]{Dablander2024}. Table~\ref{tab:hybrid_alternatives} compares these alternatives with the hybrid.

% CODE INPUT
\begin{table}[htbp]
\centering
\small
\setlength{\tabcolsep}{3pt}
\caption[Methods for combining or reducing representations.]{Comparison of methods for combining or reducing representations. For each method, the columns state whether it uses the labels and whether it keeps the original features. They also state whether one model sees features from more than one representation (Several representations) and whether adding a representation requires refitting everything already fitted (Retrain to add one). PCA = principal component analysis; PLS = partial least squares.}
\label{tab:hybrid_alternatives}
\begin{adjustbox}{max width=\textwidth}
\begin{tabular}{P{4.2cm}P{2.3cm}P{2.1cm}P{2.3cm}P{2.6cm}}
\toprule
Method & Uses the labels & Keeps original features & Several representations & Retrain to add one \\
\midrule
Concatenation \citep{Stahlschmidt2022} & no & yes & yes & yes \\
Late fusion and ensembling \citep{Steyaert2023, Valsecchi2020} & only if pooling weights are learned & yes & no & no, one more model \\
Intermediate fusion \citep{Wang2021} & yes & no & yes & yes \\
Multiple kernel learning \citep{Gonen2011, Cichonska2018} & yes & no, kernels only & yes & yes \\
PCA and PLS \citep{Jolliffe2016, Wold2001} & PCA no, PLS yes & no & yes, on a concatenation & yes \\
Flat feature selection \citep{Guyon2003} & yes & yes & yes, on a concatenation & yes \\
Group lasso \citep{Yuan2006} & yes & yes, whole groups & yes & yes \\
Adaptive hybrid representation & yes & yes & yes & search rerun, method unchanged \\
\bottomrule
\end{tabular}
\end{adjustbox}
\end{table}
\subsection{The adaptive hybrid representation}
% [notes moved to thesis_notes_moved.md]

% TODO: is this right??? I'm making shit up with NGBoost.
Here, we present KIRBy (Knowledge Integration by Representation), a framework that takes a collection of feature matrices for a given data set and outputs a customized feature set optimized for a specific prediction task. The inputs can be any type of frozen feature vector, such as fingerprints, descriptor sets, foundation model embeddings, gene expression signatures, histopathology features. Within a given domain, these individual vectors typically contain its own type of information. It's a well-established phenomenon when it comes to machine learning, there's rarely ever a one-size-fits-all approach. Take the noise robustness study above as evidence, NGBoost's predictive performance on QM9 was only ever strong with given PDV input features. However, on Caco-2, NGBoost preferred Avalon, with PDV as a close second. The model NGBoost itself produces better fits when given physicochemical properties, yet this preference can be overidden by a different prediction task. Avalon likely captures features that are pertinent to the log efflux ratio.

% TODO: REVIEW MODIFICATION
Here, we present the adaptive hybrid representation, which takes a collection of feature matrices for a given dataset and outputs a feature set customized for a specific prediction task. The inputs can be any type of frozen feature vector, such as fingerprints, descriptor sets, foundation model embeddings, gene expression signatures or histopathology features. Within a given domain, each of these vectors typically carries its own type of information. It is well established in machine learning that one approach rarely suits every problem, and the noise robustness study of Chapter~\ref{chap:robustness} provides evidence of this for NGBoost on clean labels. On QM9, NGBoost's held-out $R^2$ was highest with PDV input features (0.865, median of 10 replicates), with MHG-GNN only 0.013 lower (0.852) and ECFP4 lowest (0.706). On LogD, PDV was again highest (0.673), followed by MHG-GNN (0.658). On hERG $K_i$, MHG-GNN was highest (0.511) and PDV second (0.456), the latter being the median of the 4 of 5 folds whose clean $R^2$ was at least 0.3. On Caco-2, NGBoost appeared to prefer Avalon (0.402), with ChemBERTa at 0.404. Both numbers are medians of only 2 of 5 folds, because the other 3 fell below a clean $R^2$ of 0.3 and were excluded. Over all 5 folds, PDV is highest (0.370), ahead of Avalon (0.287) and ChemBERTa (0.256). NGBoost thus tends to produce better fits from physicochemical properties, yet a different prediction task can override this preference. Accuracy and robustness also favor different representations. On QM9 under Gaussian noise, NGBoost's highest noise robustness (\aucnorm, the normalized area under the curve of $R^2$ against noise level) was on Avalon, at 0.982 against 0.975 for PDV.

But as important as the selection of useful features is, it is also important to filter out the useless features. We could throw the kitchen sink at the models and pool together features across a wide array of representations. Not only would that pose computational challenges, but it also introduces a huge set of irrelevant features that function as noise in downstream applications (Chapter~\ref{chap:flaws}). With all of that in mind, I decided to develop a method that selects the best features for a given data set and model. In theory, it should effectively filter out noisy features and isolate the feature which carry the most relevant information.

% TODO: REVIEW MODIFICATION
Selecting useful features matters, but filtering out useless ones matters just as much. Pooling features from a wide array of representations would pose computational challenges and would also introduce a large set of irrelevant features that act as noise in downstream applications (Chapter~\ref{chap:flaws}). In an earlier exploratory study, the noise handling methods that used a distance between molecules performed worst (Section~\ref{sec:lnl_motivation}). That study served as motivation only, since it covered one dataset, the QM9 HOMO--LUMO gap, and one noise type, Gaussian noise. Chapter~\ref{chap:similarity} addresses the label-noise question over more datasets and noise types. With this in mind, I decided to develop a method that selects the best features for a given dataset and model. In principle, such a method should filter out noisy features and isolate the features that carry the most relevant information.

% TODO: review, needs manual review
The questions this chapter and Chapter~\ref{chap:similarity} address are set out in Chapter~\ref{chap:intro}.

% CODE INPUT
Formally, an adaptive hybrid representation is a feature set built by selecting up to a fixed number of features, the budget, from several base representations of the same molecules. A search chooses it for one dataset and one downstream model, so a different dataset, split or model gives a different hybrid. Because the hybrid is an input rather than a model, any standard supervised model can be trained on it. It is intended for base representations that describe different aspects of the molecule, such as PDV descriptors, ECFP4 bits, graph embeddings and sequence embeddings, from all of which the features for the task are then picked.

% CODE INPUT
The budget is split between the representations before any feature is picked. Representations differ in width by orders of magnitude, from tens to thousands of features, so in a single ranking over all their features a wide representation can crowd out a narrow one. Importance scores from different types of feature are also not on a common scale. For these reasons, each representation's features are ranked only against each other. The search also includes one pooled allocation method, the global pool, which ranks the features offered by every representation together, so that per-representation and pooled selection are compared instead of assumed. This design leaves six choices open. Three concern the features, namely how useful each feature is, how many features each representation receives and how many representations take part. The other three are how a candidate hybrid is judged during the search, the total budget, and whether the answer depends on the downstream model.

% CODE INPUT
A hybrid can also be used to measure similarity in two ways. Distances can be computed in an ordinary hybrid, or the search can pick features by how well each molecule's activity is predicted from its five nearest neighbors. Both approaches are defined in the section ``Similarity in a hybrid'' of Chapter~\ref{chap:similarity}.


\subsection{Blueprints and their chemical interpretation}
% [notes moved to thesis_notes_moved.md]

% CODE INPUT
A finished search is saved as a blueprint, which records the settings needed to rebuild the hybrid. Every selected feature traces back to its representation and to its importance rank within that representation, so the search returns information about the problem as well as a better score. The search keeps the best configuration for each pair of downstream model and budget, and one dataset therefore gives one blueprint per downstream model and budget. Comparing these blueprints at one budget shows whether tree ensembles and neural networks draw on different features.

% CODE INPUT
Some selected features can be read by a chemist. A descriptor from the PDV set carries its name, such as molecular weight or logP. A selected ECFP4 bit can be shown as the atom environment that sets it, and a MACCS key stands for a known substructure. Outside small molecules, the readable unit changes. For an antibody it is a sequence position and the residue at it, and for gene expression it is a named gene set, such as a pathway or a chromosome arm. A feature of a learned embedding, by contrast, indicates only which representation was useful and says nothing about the chemistry, for the reasons given in Chapter~\ref{chap:flaws}.

% CODE INPUT
Readable features can be used in two directions. The first is discovery, in which substructures or descriptors that recur among the kept features point to parts of the molecule that may be involved in binding or activity. Such features are a starting point for a chemist rather than a finding. The second is confirmation, in which kept features that make chemical sense for the endpoint support the claim that the selection finds real signal instead of fitting noise. Explainability is a recognized need in drug discovery \citep{JimenezLuna2020}, and inputs that can be read directly are an alternative to explaining a black-box model after the fact \citep{Rudin2019}. Two caveats apply to this reading. First, a selected feature is predictive but not necessarily causal, and because correlated features can stand in for each other, an unselected feature may matter just as much. Second, a blueprint rebuilt on other molecules can select other features, so features should be read from the fit on the dataset in question.

% CODE INPUT
A hybrid also defines a distance chosen for one task, which any method that uses nearby molecules as evidence could use in place of a standard fingerprint distance. Chapter~\ref{chap:flaws} argues for several such uses, including noise detection, nearest-neighbor prediction, applicability domain, activity-cliff detection, clustering molecules for group-aware splits and weighting predictions from separate sites. One further use is not covered there. Generative design scores each proposed molecule with a predictor \citep{Olivecrona2017}, and that predictor could take a hybrid as its input, which we leave to future work.


\subsection{Beyond small molecules}
% [notes moved to thesis_notes_moved.md]

% CODE INPUT
The search needs only a supervised task and several vector representations of the same samples, each representation being one feature matrix with one row per sample. The default search objective scores prediction only and reads no chemical structure, so the method applies to tabular data in general, that is, to any task where several encoders or measurements exist for the same inputs. What does not carry over is specific to chemistry. Reading substructures from a blueprint requires named substructure features. An optional search objective, which scores how closely distances in the hybrid track the maximum common substructure shared by pairs of molecules, applies only to molecules. The nearest-neighbor objective of Chapter~\ref{chap:similarity} needs only labels and a distance, so it does carry over and is searched on the antibody and other non-molecular datasets as well.

% CODE INPUT
One-hot sequence encodings are very sparse, because at most positions the residue never changes across a dataset, and those features therefore carry no importance. One-hot encodings can nevertheless still be useful. Among the four antibody datasets of this chapter, one-hot is a base representation on nanobody thermostability only, while the position-wise blocks substitution matrix (BLOSUM) encoding is a base representation on humanness and trastuzumab--HER2 binding. One-hot was among the selected representations on \textbf{[RESULT NEEDED: the antibody datasets where one-hot was among the selected representations, with the score in the dataset's metric]}. We recommend language-model sequence embeddings over raw one-hot encodings.

% CODE INPUT
After small molecules, antibodies received the most attention in this work. This chapter uses four antibody datasets, Koenig expression, nanobody thermostability, humanness and trastuzumab--HER2 binding (Table~\ref{tab:datasets}). Their base representations are mostly embeddings from general protein language models and from antibody-specific language models, some of which cover only the complementarity-determining region (CDR) loops or come from early or middle layers of ESM-2 (Table~\ref{tab:reps_antibodies}). Chemical reactions serve as a proof of concept on two Buchwald--Hartwig yield datasets. Their base representations are computed on the reactants, on the products minus the reactants, or on the whole reaction, and product fingerprints enter only as augmentation (Table~\ref{tab:reps_reactions}). Nucleotide sequences are a second proof of concept, on the H3K4me3 histone-mark and TATA promoter tasks (Table~\ref{tab:reps_nucleotides}). In cancer genomics, on data from The Cancer Genome Atlas (TCGA), each representation is a different measurement of the patient, namely gene expression, methylation (lung task only), copy number, microRNA (miRNA), protein and slide images (Table~\ref{tab:reps_tcga}). In the other domains, each representation is a different encoder of one input, and the search treats both cases in the same way.

% CODE INPUT
Two conditions must hold on a dataset for the hybrid to help. The first is a sufficient number of samples. \textbf{[RESULT NEEDED: the number of samples below which the hybrid stops beating the best single representation; about 200 samples is expected]}. The second is base representations whose features carry importance, since where every representation already scores poorly, there are too few useful features to improve on.


\section{Methods: the hybrid search}
% [notes moved to thesis_notes_moved.md]

% CODE INPUT
Figure~\ref{fig:hybrid_search_space} gives an overview of how a hybrid is built from several base representations of the same samples.

% CODE INPUT: figure added from ~/Downloads/Pruning.png (30 Sep). ~/Downloads/Hybrid Search Space.png (8 Jul) is the
% same diagram with "Dataset" above the inputs. TODO: review -- pick one.
\begin{figure}[htbp]
    \centering
    \includegraphics[width=\textwidth]{figures_code_input/hybrid_search_space.png}
    \caption[Schematic overview of the hybrid search.]{Schematic overview of the hybrid search. The inputs (left) are the base representations, the labels and the target metric, which are split into five cross-validation folds, and the single representations are scored as baselines. Within each fold, the search covers importance methods, allocation methods, budgets, internal evaluators and downstream models (center). Adaptive pruning, when switched on, uses the first two folds to drop settings. The best hybrid is then selected for each downstream model, and its settings are stored as a blueprint (right).}
    \label{fig:hybrid_search_space}
\end{figure} The search is a grid in which every combination of setting values is built and scored inside the cross-validation folds. When adaptive pruning is switched on, it drops some combinations after the first two folds. A winner is kept for each pair of downstream model and budget, where the budget is a maximum number of features and a hybrid can have fewer. The comparison we report runs the whole search inside each outer training fold and scores the winner once on the untouched outer test fold.

% CODE INPUT
The subsections below cover the steps from base representations to a hybrid, each searched setting and its rationale, augmentation with small or specialized feature sets, and two ways of making the search cheaper.


\subsection{From base representations to a hybrid}
% [notes moved to thesis_notes_moved.md]

% TODO: review
% TODO: review, needs manual review
Given $K$ frozen representation matrices $\mathbf{R}_1, \dots, \mathbf{R}_K$ for the same $n$ samples, a label vector $\mathbf{y}$, and optional group assignments for cross-validation, the search finds which features from which sources best predict for a specified metric and downstream model. The output is a hybrid feature matrix of at most $b$ features, where $b$ is the feature budget, plus any augmentation features added outside it. The search also returns a blueprint that records all parameters needed to reproduce the identical feature selection on new data. The search first scores each source independently, then scores feature importance on training data within each cross-validation fold, allocates the budget among sources according to baseline performance, and scores candidate hybrid configurations via cross-validation. We retain the best configuration for each combination of downstream model and budget. Every selected feature traces back to its source and importance rank.

% TODO: REVIEW MODIFICATION
The search is given $K$ frozen representation matrices $\mathbf{X}^{(k)} \in \mathbb{R}^{n\times d_k}$, $k = 1, \dots, K$, for the same $n$ samples, whose widths $d_k$ range from tens to thousands of features. It is also given a label vector $\mathbf{y}$ and one group label per sample for cross-validation, for example its scaffold. It further needs the dataset's metric, the task type, one or more downstream models and a list of budgets. From these inputs, it finds which features from which representations best predict for that metric and downstream model. The output is a hybrid feature matrix of at most $B$ features, where the feature budget $B$ is a maximum rather than a target. Uncapped augmentation adds its extra features outside the budget, whereas capped augmentation stays within it. The output also includes a blueprint that records the settings needed to rebuild the hybrid. The search first scores each representation independently and then scores feature importance on training data within each cross-validation fold. It then allocates the budget among representations. Four of the five allocation methods do this according to each representation's standalone performance, while the global pool instead ranks the offered features from all representations together. Candidate hybrid configurations are scored by cross-validation, and we retain the best configuration for each combination of downstream model and budget. Every selected feature traces back to its representation and importance rank.

% TODO: review, needs manual review
Inputs from different sources differ in scale, distribution, and missingness patterns, and naively pooling them before splitting would leak validation-set statistics into the training procedure, a well-documented source of optimistic bias in feature selection studies \citep{Ambroise2002, Vabalas2019}. All preprocessing therefore runs within each cross-validation fold. We impute missing values with per-feature medians computed on the training fold and drop features that are entirely missing. Every feature of the hybrid, binary fingerprint bits included, is standardized; only when a single representation is scored on its own are binary features left as 0/1. We standardize to zero mean and unit variance using training-fold statistics only.

% TODO: REVIEW MODIFICATION
Inputs from different representations differ in scale, distribution and missingness patterns. Pooling them naively before splitting would leak validation-set statistics into the training procedure, a well-documented source of optimistic bias in feature selection studies \citep{Ambroise2002, Vabalas2019}. All preprocessing therefore runs within each cross-validation fold. When a single representation is scored, we impute missing values with per-feature medians computed on the training fold and drop features that are entirely missing. When features are selected for the hybrid, a missing value is instead set to 0 before standardization. Every feature of the hybrid, binary fingerprint bits included, is standardized to zero mean and unit variance using training-fold statistics only. Binary features are left as 0/1 in one case only, when the search scores a single representation to pick the comparator, whereas a single representation scored on the outer test fold is standardized like the hybrid. A representation wider than 5,000 features is first reduced, inside each fold, to the 500 features on which a ranking model splits most often. It keeps 200 features instead when the search is given fewer than 1,000 molecules. The ranking model is a LightGBM model of 50 trees of depth at most 5, fitted on the fold's training molecules. It ranks features by how many times its trees split on them rather than by gain. This reduction applies only when the search scores a single representation on its own. For such a representation, the comparator is therefore picked on its top 500 (or 200) features and then scored on the outer test fold at full width, while the hybrid search always sees the full width.

% CODE INPUT
Each dataset is split into outer folds, 5 by default. The folds are group-aware, so that no group appears on both sides of a split, and classification folds are also stratified by class. A representation learned from the labels is retrained inside each outer training fold, on its training molecules only. The search then runs its own group-aware cross-validation, with 5 folds by default, on the molecules it is given. Inside each training fold of the search, one group-aware internal split holds out about 20\% of the molecules as whole groups. This split is the first fold of a 5-fold grouped split and is therefore deterministic. When a training fold holds fewer than five groups, both splits use as many folds as there are groups. For example, Koenig expression has five groups, so each outer training fold holds four. The search inside that fold uses 4 folds of three groups each, and the internal split holds out one of the three. Importance and each representation's standalone score come from this training fold only, and the outer test molecules score only the finished hybrid.

% CODE INPUT
For each representation $k$, the importance method gives every feature one score $s^{(k)}_j \ge 0$, and these scores rank features within that representation only. Importance is computed on the internal training part of the fold. The standalone score $r_k$ comes from the internal evaluator, fitted on representation $k$ alone and scored on the internal held-out part in the dataset's metric. We rescale $r_k$ to $\tilde r_k$ so that 0 means ``no better than a trivial predictor'', which makes the budget split independent of the dataset's metric. For Spearman, Pearson, the Matthews correlation coefficient (MCC) and $R^2$, $\tilde r_k = \max(0, r_k)$. For AUROC and the concordance index, chance is 0.5, so $\tilde r_k = \max(0, 2(r_k - 0.5))$. For RMSE, $\tilde r_k = \max(0, 1 - \mathrm{RMSE}_k/\mathrm{RMSE}_0)$, where $\mathrm{RMSE}_0$ is the error of predicting the training mean, and the mean absolute error is rescaled in the same way with the training median as the trivial predictor. For accuracy and weighted F1, $\tilde r_k = \max(0, (r_k - r_0)/(1 - r_0))$, where $r_0$ is the score of always predicting the majority training class. The trivial predictor is scored on the same held-out molecules. A representation whose fit fails receives weight 0 and is recorded as a failure.

% CODE INPUT
An optional representation count $N$ keeps only the $N$ representations with the highest standalone scores on this fold. Each remaining representation receives a weight $w_k = g(\tilde r_k)$, where the function $g$ is set by the allocation method. The exact share of representation $k$ is $b^*_k = B\, w_k / \sum_l w_l$ features, and we round the shares by largest remainder. Each representation first receives $\min(\lfloor b^*_k \rfloor, d_k)$ features, and the leftover features are then given out one at a time in order of largest fractional part, skipping representations already at full width. The result satisfies $\sum_k b_k = \min(B, \sum_k d_k)$ and $b_k \le d_k$, and a representation with $w_k > 0$ keeps at least one feature. If every weight is 0, the budget is split equally and the run reports this, and if no representation can be scored at all, the configuration is refused. The global pool does not use this rounding. Instead, each representation offers candidate features, and one joint ranking keeps at most $B$ of them.

% CODE INPUT
Let $S_k$ be the top $b_k$ features of representation $k$ by $s^{(k)}_j$, among those with $s^{(k)}_j > 0$. The same cut at 0 applies to the features offered to the global pool and to its final ranking. If no feature of a representation scores above 0, its full $b_k$ is taken. The cut exists because tree importance gives exactly 0 to every feature on which no tree split, and filling the budget from those features would pick them by their position in the matrix. The hybrid is $\mathbf{H} = [\mathbf{X}^{(1)}_{:,S_1}, \dots, \mathbf{X}^{(K)}_{:,S_K}]$, with $\sum_k |S_k| \le B$, and it can therefore be narrower than its budget. The matrix $\mathbf{H}$ is then preprocessed inside the fold, as described above.

% CODE INPUT
A configuration assigns one value to every searched setting. For each configuration, we fit the downstream model on the hybrid's training molecules in each fold of the search, score it on that fold's held-out molecules and average over the folds. The winner is picked separately for each pair of downstream model and budget and never across models. Because that score was used to choose among many configurations, it is optimistic \citep{Varma2006, Cawley2010}, and the comparison we report is therefore nested. The whole search runs inside each outer training fold, the winning settings are rebuilt on that fold, and the result is scored once on the untouched outer test fold. Both the winner at each budget and each downstream model's best across budgets are scored in this way. The inner search score exceeds the outer test score by \textbf{[RESULT NEEDED: the gap between inner search score and outer test score, per dataset, in each dataset's own metric]}.

% CODE INPUT
Chapter~\ref{chap:similarity} uses two variants of this search. Under label noise, every quantity computed on the training molecules uses the noisy labels, and held-out scoring uses clean labels unless stated otherwise. Under the nearest-neighbor objective, each molecule's activity is predicted from its five nearest neighbors in the feature space. This score replaces both the standalone score and the downstream-model score, so no downstream model or internal evaluator is used.

% TODO: Figure 5.1: the search pipeline schematic. Does not exist yet.


\paragraph{Blueprints.}

% TODO: review, needs manual review
To be able to better examine the selected features and to regenerate them on new datasets, we produce blueprints. After a hybrid feature selection search is completed, its results are stored as a blueprint containing the ``recipe'' required to recreate this exact hybrid representation on any data. These blueprints store the selected picks for the pre-specified budget, the full list of input representations searched and how many of them the recipe used, feature importance scoring method, allocation strategy, and internal evaluator. It contains other runtime information like the step size, number of folds, the seed, the metric being used to score, the task type (regression or classification), whether augmented features were being used, and the per-representation scores.

% TODO: REVIEW MODIFICATION
To examine the selected features more closely and to rebuild the hybrid on new datasets, we produce blueprints. A blueprint is the recipe for the winning hybrid, that is, the search settings needed to rebuild it. Every scored configuration has one, and the search returns the winner's blueprint with its result. On other molecules, a rebuild re-scores the representations and re-picks the features, so it need not recreate the same hybrid. A blueprint stores the budget, the full list of input representations searched and how many of them the recipe used. It also stores the feature importance method, allocation method, internal evaluator and objective. Further runtime information includes the number of folds, the seed, the metric used for scoring and its direction, and the task type (regression or classification). The augmentation settings are stored as well, covering how many extra features to take, from which candidates, and the rule for admitting one. Some information is recorded for later analysis only. This comprises each representation's best score across the downstream models, the number of features each representation received in each fold, and the augmentation features kept in each fold. These records show which representations fed the hybrid and whether this holds across folds. The stored score differs from the standalone score that splits the budget, because it is measured on every molecule the search saw, and a rebuild therefore never reads it.

% TODO: REVIEW MODIFICATION -- blueprint shown as a data structure. Every field name and type is copied from the
% Blueprint class in src/kirby/find_best_hybrid.py:1415-1553 (read 2026-09-30). The comments on which fields a rebuild
% reads come from build_fold_hybrid in tests/experiments/_shared.py:5091 (the per-fold rebuild the experiments use).
% That rebuild reads rep_scores in one case (see the comment on that field). The paragraph above says a rebuild never
% reads the stored score; that sentence needs your decision.
Listing~\ref{lst:blueprint} gives every field of a blueprint together with its type.
\begin{lstlisting}[float=htbp, breaklines=true, basicstyle=\scriptsize\ttfamily, columns=fullflexible, keepspaces=true, frame=single, caption={Fields of a blueprint and their types. Comment lines group the fields, and each field is given with its type and, where needed, a comment on its meaning.}, label=lst:blueprint, captionpos=b]
Blueprint
  # What the search chose. A rebuild reads all of these.
  reps                  : list[str]           # every representation searched, not only those used
  rep_count             : int or None         # how many of them may contribute; None = all
  budget                : int                 # most features the hybrid may have
  allocation_method     : str
  importance_method     : str
  internal_evaluator    : str
  objective             : str                 # 'predictive' or 'similarity'

  # How the search scored configurations.
  metric                : str
  lower_is_better       : bool
  task_type             : str                 # 'regression' or 'classification'
  n_folds               : int
  seed                  : int

  # Augmentation settings. A rebuild reads all of these.
  augmentation_strategy : str
  aug_per_source_budget : int, or {source: int}, or None
  augmentation_margin   : float
  augmentation_fraction : float or None
  candidate_ranking     : str
  augmentation_threshold: float or None
  augmentation_budget   : int
  aug_max_candidates    : int

  # Stored for inspection.
  rep_scores            : {representation: float}
                          # each representation's best score over every molecule the search saw;
                          # the rebuild uses it to choose which representations to keep only when
                          # rep_count is set and no per-fold scores are given
  allocation_per_fold   : {fold: {representation: int}}      # features each representation got
  kept_augmentations    : {fold: {source: list[int]}}         # extra features kept, by index

  # Left from a removed method.
  step_size             : int                 # used only by greedy forward selection
\end{lstlisting}


The blueprint does not store selected feature indices, those can only be recovered when the blueprint is refit on the data, which re-runs the feature selection process deterministically from that recipe and seed.

% TODO: REVIEW MODIFICATION
The blueprint does not store the selected base-representation features, although it does store how many features each representation received in each fold and which augmentation features were kept. A blueprint can be reused in two ways. The first re-runs the recipe on new labeled data and recomputes importance and standalone scores, so that the features adapt to the new dataset. This costs a new selection and can give a different feature set. The second takes the exact features of a labeled fit from the new molecules. This requires no labels and is therefore the only option for unlabeled molecules, but it assumes that the same features matter on the new data. Re-running the recipe without labels is refused, because every importance score would be 0 and the selection would fall back to feature order. Which of the two approaches transfers better to a smaller related dataset is \textbf{[RESULT NEEDED: re-run recipe vs. reused exact features, large to small related dataset]}. Importance and allocation can also be computed separately within k-means clusters of the molecules' ECFP fingerprints, to ask whether the budget split changes across chemical space.

It also provides a registry of selected features, stating where that feature came from. On some representations, we added an extension where features are mapped back to their function. For instance, ECFP4 and MACCS bits are mapped to SMILES arbitrary target specification (SMARTS) substructures. This blueprint will tell you which of those substructures were selected. This is not possible for learned embeddings for which the blueprint will return the bare representation name and dimension index. This mapping provides a layer of explainability. It can be used both to discover what features are useful for your task and to confirm that the features being selected are in fact lining up with the underlying known chemistry.

% TODO: REVIEW MODIFICATION
A separate explanation step, run on a finished hybrid, provides a registry of the selected features that states where each feature came from, since the blueprint itself does not store feature names. For some representations, we added an extension that maps features back to their function. An ECFP4 bit is named by the atom environment that sets it, written as a SMILES fragment rather than a SMILES arbitrary target specification (SMARTS) pattern. The fragment is taken from the first molecule, of at most 50 supplied, that has the bit set. MACCS keys are named by their SMARTS pattern and a short name, and the features of the descriptor vector carry the descriptor names. The explanation step thus reports which of these substructures and descriptors were selected. Outside small molecules, protein composition, descriptor and k-mer features are named too, as are antibody composition, one-hot and BLOSUM positions, structural and inverse-folding scores, gene-expression pathway sets and PCA components. MAP4 and DRFP features are returned as a bit index only. Avalon, Mordred, atom-pair and graph-kernel features, and every learned embedding, return only the bare representation name and dimension index. This mapping adds a layer of explainability. It can be used to discover which features are useful for a task and to confirm that the selected features agree with known chemistry.

\subsection{The search space}
% [notes moved to thesis_notes_moved.md]

% TODO: REVIEW MODIFICATION
When I began to implement this idea, I expected a fairly straightforward way to select the most informative features, but several questions soon arose. How should the importance of a feature be measured? Given a handful of input representations, how many features should be taken from each? How should a candidate feature set be scored? As with the choice of representation, different prediction tasks preferred different answers to each of these questions. The search therefore treats each question as a setting and chooses its value per dataset. These settings are the importance method, the allocation method together with the budget and the representation count, and the internal evaluator.

% CODE INPUT
Every combination of setting values in the grid is built and scored, and nothing is fixed in advance except the values each setting may take. The claim of this chapter is that no single setting suits every dataset, and a grid measures this claim instead of assuming it. The motivation is the familiar result that no learning method is best on every problem \citep{Wolpert1996}, although that result does not prove the claim for hybrids. We search the full product of values, instead of one setting at a time, because the settings interact. For example, the best budget depends on the downstream model, as the nanobody thermostability example below shows.

% CODE INPUT
Six settings are searched, namely the downstream model, the budget, the representation count, the importance method, the allocation method and the internal evaluator. Their values at the full tier and at the fast tier are listed in Table~\ref{tab:search_tiers}. The defaults and the number of configurations for each downstream model are given in Additional files, Section~\ref{app:searchspace}. The importance method is LightGBM gain, TreeSHAP or DeepLIFT. It is searched because importance depends on the model as well as on the data \citep{Dong2020}. The three methods come from different model families and rank features differently, and a different ranking builds a different hybrid. The allocation method decides how the budget is split among the participating representations, and five allocation methods are searched. The internal evaluator scores each representation on its own. It is the matched evaluator (a model of the downstream model's own family), random forest, LightGBM or a deep neural network (DNN). Two further choices are not searched within a single search. The objective is prediction by default, which Chapter~\ref{chap:similarity} replaces with predicting each molecule's label from its nearest neighbors in the hybrid. Augmentation is absent, uncapped or capped. Each objective and each augmentation condition is therefore a separate search.

% CODE INPUT
The downstream model is the model trained on the finished hybrid. Five downstream models are searched, namely random forest, XGBoost, LightGBM, SVM and DNN, and a Gaussian process is added on datasets under 1,000 molecules. Three of the 16 datasets are that small, but run time rules the Gaussian process out on two of them, so it runs on humanness only. The downstream model is searched because it determines what counts as a useful feature and how many features it can use. The best hybrid for one model therefore need not be the best for another.

% CODE INPUT
The budget is the maximum number of features the hybrid may have, and eight budgets are searched (10, 25, 50, 100, 250, 500, 1,000 and 2,000 features). The budget is searched because too few features leave informative ones out, whereas too many add uninformative features that, on small datasets, add room to overfit. Budgets above the number of training molecules are kept. The best budget depends on the downstream model. In a preliminary run on nanobody thermostability with 5 folds, SVM performed best at 300 features, with a Spearman correlation of 0.6384. DNN, LightGBM and random forest performed best at 50 features, and XGBoost at 100. That run used an earlier list of budgets, which included 300 features. \textbf{[RESULT NEEDED: the best budget for each downstream model on nanobody thermostability, from the full runs]}

% CODE INPUT
Each dataset has a fixed set of base representations, chosen before the search (Table~\ref{tab:dataset_reps}). Fourteen of the 16 datasets have 10, H3K4me3 has 8 and TATA promoter has 9. The representation count, which takes the values 2, 3, 5 or all, sets how many of them may contribute features to the hybrid. On each fold, the representations with the highest standalone scores are retained. This is the only step that drops whole representations before allocation, and correlation between representations plays no part in it. The count is searched because datasets differ in how many representations are worth using, and pooled feature selection over all representations has no such setting. Every allocation method is paired with every count. In preliminary runs, some datasets performed best with 2 or 3 representations and others with all 10. \textbf{[RESULT NEEDED: which datasets do best with 2 or 3 representations and which with all of them, from the full runs]}

\subsubsection{Feature importance scoring}
% [notes moved to thesis_notes_moved.md]

% TODO: this is out of date

% What were the qualifications for correlation/filtering allocation strategies? Can I formalize this a bit?
% TODO: review, needs manual review
Not all features from a given base representation are equally informative, and there are a variety of techniques that can assess each feature's relative importance. Tree-based feature importance measures favor features with high split gain, while neural attributions such as DeepLIFT measure how much each feature moves a neural network's output away from its output on a reference input \citep{Strobl2007, Ancona2018}.
We tested whether, for the purposes of this task, there was a single feature importance method that worked across datasets and model architectures. Not only was there no clear winner, but there was evidence that certain networks preferred tree-based measures of importance while certain tree ensembles preferred gradient-based approaches. We landed on a set of three whose predictions did not correlate strongly with each other and whose selected features performed well as model input.

% TODO: REVIEW MODIFICATION
Not all features from a given base representation are equally informative, and a variety of techniques can assess each feature's relative importance. The search ranks the features of each representation before choosing any, using an importance method that gives every feature of one representation a single non-negative score. Each representation receives its own freshly fitted model on its own features only, and the downstream model is not used. Only the order within a representation is used, and because these scores are never compared across representations or across methods, no normalization is applied. Scores are computed inside each fold on the internal training part only, which consists of 80\% of that fold's training molecules, split by group. Before fitting, missing values are set to 0 and each feature is standardized to mean 0 and variance 1 on the same molecules. Scores are computed once per fold and method and are then reused by every allocation method, budget, internal evaluator and downstream model. Split-based tree importance is biased toward features with many possible split points \citep{Strobl2007}. This bias matters here, because binary fingerprint bits sit beside continuous embeddings. Neural attributions such as DeepLIFT instead measure how much each feature moves a neural network's output away from its output on a reference input \citep{Ancona2018}.

% TODO: REVIEW MODIFICATION
We tested whether a single feature importance method worked across datasets and downstream models for this task. We compared the methods' feature rankings on 7 datasets and the hybrids built from those rankings on one dataset with two downstream models. We also counted which method gave the best configuration on 5 datasets with 5 downstream models. There was no clear winner, and the preferred method changed with the dataset. With a DNN downstream model, LightGBM gain was preferred on four of five datasets, whereas on LogD alone the tree ensembles XGBoost and LightGBM preferred DeepLIFT. We settled on a set of three methods whose rankings did not correlate strongly with each other and whose selected features performed well as model input.

% TODO: background on each of these, also how do I name these in text
% TODO: review, needs manual review
We score features using LightGBM gain, DeepLIFT, and TreeSHAP. LightGBM gain is a tree based method which.... TreeSHAP is explanation-based, meaning it... DeepLIFT relies on neural attribution which...

% TODO: REVIEW MODIFICATION
We score features using LightGBM gain importance, SHapley Additive exPlanations (SHAP) values for tree ensembles (TreeSHAP) and DeepLIFT (Deep Learning Important FeaTures) \citep{Shrikumar2017}, which come from three families, namely split gain, Shapley values and reference-based attribution. LightGBM gain measures how much a feature reduced the training loss where trees split on it \citep{ke2017lightgbm}. Its model consists of 50 gradient-boosted trees of maximum depth 6, and a feature's score is the loss reduction of every split on it, summed over all trees, or 0 if no tree splits on it. TreeSHAP measures how much a feature moves each molecule's prediction, as an exact Shapley value for tree models \citep{Lundberg2020}. Its model is a random forest of 50 trees with maximum depth 10 \citep{Breiman2001}, and a feature's score is its mean absolute Shapley value over the training molecules. DeepLIFT measures how much a feature moves the output of a small network away from its output on a reference input \citep{Shrikumar2017, Kokhlikyan2020}. The network has two hidden layers of 64 and 32 ReLU units and is trained for 100 full-batch epochs of Adam at a learning rate of 0.001 \citep{Kingma2015}. A feature's score is its mean absolute attribution over the training molecules under DeepLIFT's Rescale rule. The reference input is all zeros, which after standardization places every feature at its training mean. For a fingerprint bit, this mean is the fraction of training molecules that have the bit set. On binary labels, TreeSHAP and DeepLIFT use the positive class, and on multiclass labels they average absolute values over the classes. For LightGBM gain and TreeSHAP, a representation wider than 2,000 features is pre-screened. Only its 2,000 highest-variance features, with variance measured before standardization, are fitted and the rest score 0, whereas DeepLIFT scores every feature. On Caco-2, the median share of features scoring above 0 was 30\% for LightGBM gain, 70\% for TreeSHAP and 90\% for DeepLIFT. This preliminary median is taken over 10 representations and 3 folds. Table~\ref{tab:importance_methods} summarizes the three importance methods.

% CODE INPUT
\begin{table}[htbp]
\centering
\caption[Settings of the three importance methods searched.]{Settings of the three importance methods searched. Each method is fitted separately on every representation, inside each fold, on the internal training part. For each method, the columns give its family, the model fitted per representation, how the score of one feature is computed, how classification labels are handled and how representations wider than 2{,}000 features are treated. Settings not listed are the library's defaults. DeepLIFT = Deep Learning Important FeaTures; ReLU = rectified linear unit.}
\label{tab:importance_methods}
\small
\begin{adjustbox}{max width=\textwidth}
\begin{tabular}{@{}P{0.15\linewidth}P{0.24\linewidth}P{0.22\linewidth}P{0.17\linewidth}P{0.16\linewidth}@{}}
\toprule
Method (family) & Model fitted per representation & Score of one feature & Classification labels & Representations over 2{,}000 features \\
\midrule
LightGBM gain (split gain) & LightGBM, 50 trees, maximum depth 6 & Loss reduction of every split on the feature, summed over all trees; 0 if never split on & One model on the class labels & Only the 2{,}000 highest-variance features are fitted; the rest score 0 \\
TreeSHAP (Shapley values) & Random forest, 50 trees, maximum depth 10 & Mean absolute Shapley value over the training molecules & Binary: positive class. Multiclass: mean absolute value over classes & As for LightGBM gain \\
DeepLIFT (reference-based attribution) & Small neural network, ReLU hidden layers of 64 and 32 units, 100 full-batch epochs of Adam, learning rate 0.001 & Mean absolute attribution against an all-zero reference, Rescale rule & Binary: positive class. Multiclass: mean absolute value over classes & Every feature is scored \\
\bottomrule
\end{tabular}
\end{adjustbox}
\end{table}

When searching for the features to use in the hybrid, thse feature importance scoring methods are swept.

% TODO: REVIEW MODIFICATION
All three feature importance methods are swept in the full search, whereas the fast tier uses LightGBM gain only (Table~\ref{tab:search_tiers}). Adaptive pruning is off unless a search requests it. When it is on, importance methods that rarely give the best configurations on the first two folds are dropped from the remaining folds, although at least one method is always kept. Pruning is described in the subsection ``Fast tier and adaptive pruning''.

% CODE INPUT
Two comparisons were run, the first of how the methods rank features and the second of the hybrids those rankings build. The ranking comparison scored 8 importance methods on 7 datasets, namely LogD, hERG $K_i$, Buchwald--Hartwig yield, nanobody thermostability, antibody FcRn retention, breast invasive carcinoma (BRCA) PAM50 subtype and lung adenocarcinoma (LUAD) vs lung squamous cell carcinoma (LUSC). For each pair of methods, it measured the Spearman correlation between their scores over every feature of every representation. The three retained methods rank features differently. Across the 7 datasets, the correlation ranged from 0.20 to 0.64 between LightGBM gain and TreeSHAP, from 0.21 to 0.67 between LightGBM gain and DeepLIFT, and from 0.56 to 0.99 between TreeSHAP and DeepLIFT. These numbers are preliminary, because each method was fitted once on the whole dataset rather than inside folds.

% CODE INPUT
The hybrid comparison used Caco-2, with 2,158 molecules, 10 representations, 3 scaffold folds and a budget of 100 features. Hybrids were built from each method's ranking and scored with random forest and LightGBM downstream models. Three control rankings were added, namely feature variance, Laplacian score and a random order. The measure was the RMSE of the best single representation minus the RMSE of the hybrid, in log$_{10}$ efflux ratio units, so that a positive value means the hybrid has the lower error. The preliminary results were $+0.0144$ for DeepLIFT, $+0.0100$ for LightGBM gain, $+0.0067$ for TreeSHAP and $-0.0118$ for the random ranking. The gap between DeepLIFT and LightGBM gain, 0.0044 RMSE, is smaller than the fold-to-fold spread. The methods also differ at the top of their rankings. Of each representation's 100 top-ranked features on Caco-2, the retained pairs of methods share a median of 28 to 54 features.

% CODE INPUT
Three rules determined which methods were retained. First, a method must score every feature of every representation for regression, binary and multiclass labels, which all three retained methods do. Second, of two methods whose rankings agree above a Spearman correlation of 0.95, we keep the faster one. Under this rule, DeepLIFT replaced integrated gradients, which agreed with it above 0.95 on 4 of 6 datasets, and gradient SHAP, which agreed above 0.95 on 3 of 6. Third, we keep methods that rank features differently, because a different ranking builds a different hybrid. TreeSHAP came last of the three on Caco-2 and was kept under this third rule. Only two filters act on features before allocation, the cut at a score of 0 and the pre-screen on representations wider than 2,000 features. On a 2,048-bit fingerprint, the pre-screen drops the 48 lowest-variance bits. On hERG $K_i$ all 48 are among the 105 constant ECFP4 bits, while on QM9, which has no constant ECFP4 bits, they are near-constant ones. Near-constant features otherwise remain candidates and are ranked by importance.

% CODE INPUT
No single importance method suits every dataset and downstream model. A preliminary run used 5 datasets (nanobody thermostability, hERG $K_i$, humanness, LogD and BRCA PAM50 subtype), 5 folds, 10 representations, 5 downstream models and budgets of 25, 50, 100, 150 and 300 features. For each of the 25 pairs of downstream model and budget per dataset, we counted which importance method gave the best configuration, and Table~\ref{tab:importance_wins} gives the counts. LightGBM gain won the most pairs on three datasets and tied with TreeSHAP on humanness. DeepLIFT won the most on one dataset, LogD, with 17 of its 25 pairs, and outside LogD it gave the best configuration for 2 of 100 pairs. With the DNN downstream model, LightGBM gain won 17 of 20 budgets outside LogD, whereas on LogD the DNN took DeepLIFT at 3 of 5 budgets. On LogD, XGBoost and LightGBM each took DeepLIFT at 5 of 5 budgets, and random forest took LightGBM gain at 5 of 5. \textbf{[RESULT NEEDED: the same win counts, per dataset and importance method, from the full runs]}

% CODE INPUT
\begin{table}[htbp]
\centering
\caption[Importance method of the best configuration in a preliminary run.]{Importance method of the best configuration in a preliminary run on five datasets. A win is one pair of downstream model and budget at which that importance method gave the best configuration, out of 25 pairs per dataset. The best score is the best configuration reached with that method, in the dataset's own metric. The run used 5 folds, 10 representations, 5 downstream models and budgets of 25, 50, 100, 150 and 300 features. Lower is better for the root mean squared error (RMSE), and higher is better for Spearman correlation, the area under the receiver operating characteristic curve (AUROC) and weighted F1.}
\label{tab:importance_wins}
\small
\begin{adjustbox}{max width=\textwidth}
\begin{tabular}{@{}llcccccc@{}}
\toprule
 & & \multicolumn{3}{c}{Wins (of 25)} & \multicolumn{3}{c}{Best score} \\
\cmidrule(lr){3-5}\cmidrule(lr){6-8}
Dataset & Metric & \makecell{LightGBM\\gain} & TreeSHAP & DeepLIFT & \makecell{LightGBM\\gain} & TreeSHAP & DeepLIFT \\
\midrule
\makecell[l]{Nanobody\\thermostability} & Spearman & 15 & 10 & 0 & 0.6696 & 0.6471 & 0.6440 \\
hERG $K_i$ & RMSE (p$K_i$) & 14 & 11 & 0 & 0.5589 & 0.5605 & 0.5728 \\
Humanness & AUROC & 12 & 12 & 1 & 0.9723 & 0.9713 & 0.9689 \\
LogD & RMSE (log units) & 8 & 0 & 17 & 0.4250 & 0.4345 & 0.4170 \\
BRCA PAM50 subtype & weighted F1 & 21 & 3 & 1 & 0.9097 & 0.9018 & 0.9038 \\
\bottomrule
\end{tabular}
\end{adjustbox}
\end{table}


\subsubsection{Feature allocation methods}
% [notes moved to thesis_notes_moved.md]

% TODO: allocation strategies completely out of date
% Check framework doc?
% TODO: review, needs manual review
The allocation strategy divides the total feature budget among sources, and we search five allocation strategies. Four divide the budget in fixed shares from transformed standalone scores (linear, squared, softmax, rank decay). The fifth, a global pool, lets features from all sources compete in one importance ranking. Once we have the allocation, we select the top features from each source by importance score and concatenate them to form the hybrid.

% TODO: REVIEW MODIFICATION
The allocation method divides the total feature budget among the representations. It decides how many of the budget's $B$ features each representation contributes, while the importance method decides which ones. In total, we search five allocation methods. Four give each representation a fixed share of the budget, computed from its rescaled standalone score (linear, squared, softmax and rank-decay allocation). The fifth, the global pool, lets the candidate features offered by every representation compete in one importance ranking, in which only the offered candidates take part. Under the four fixed-share methods, we select the top features from each representation by importance score and join them to form the hybrid.

% TODO: completely replace this is SO out of date, I need a version of this for all the allocation strategies I'm using, as well as the knobs

% \begin{enumerate}
%     \item \textbf{Split:} Divide training data into allocation-train (80\%) and validation (20\%) subsets using scaffold splitting.

%     \item \textbf{Initialize:} Set current feature set to empty, remaining budget to $B$.

%     \item \textbf{Rank features:} For each representation $r \in \{base_1, base_2, \dots base_m\}$, train RFs using the base representations individually, then rank features by mean decrease in impurity.

%     \item \textbf{Iterate:} While remaining budget $\geq$ step\_size and patience not exhausted:
%     \begin{enumerate}
%         \item For each representation, try adding the top step\_size candidates from its ordered list to the current feature set.
%         \item Train a RF using the tentative feature set.
%         \item Evaluate $R^2$ on the validation set.
%         \item Select the representation whose addition yields the highest validation $R^2$.
%         \item Permanently add step\_size features from the winning representation; remove from its candidate list.
%         \item If validation $R^2$ does not improve for patience consecutive rounds, terminate early.
%     \end{enumerate}

%     \item Re-compute importance scores for the features of the new hybrid on the full training data. For each representation, select the top $n_r$ features according to the allocation determined above, where $n_r$ is the number of features allocated to representation $r$ during iteration.
% \end{enumerate}

% This greedy approach works better on larger datasets, as the size of the validation set may not be enough to gain meaningful insights from.

% TODO: this is not what's happening
% TODO: review, needs manual review
The four fixed-share methods are variants of performance-weighted allocation, where the distribution of features from each base representation is proportional to that representation's standalone predictive performance on the training set. This works under the assumption that representations with stronger predictive performance contribute more valuable features. It is worth noting that this is not always the case, as high-quality, large representations may rely on a small subset of valuable features.

% TODO: REVIEW MODIFICATION
The four fixed-share methods set each representation's share from its standalone score. This score is measured on the internal validation part of the training fold rather than on the molecules on which the internal evaluator was fitted. Only linear allocation makes the share proportional to that score, while the other three transform it first. All four assume that representations with stronger predictive performance contribute more valuable features.

% TODO: review, needs manual review
\begin{enumerate}
    \item Fit the internal evaluator on each base representation independently, on the internal training part of the fold.
    \item Score each representation on the internal validation part in the dataset's own metric, rescaled so 0 means no better than a trivial predictor.
    \item Allocate features proportionally:
    \[
    b^*_k = B \, \frac{w_k}{\sum_{l} w_l}, \quad w_k = g(\tilde r_k) % TODO: REVIEW MODIFICATION -- notation changed from n_r, w_r to b_k, w_k to match the rest of the chapter
    \]
    \item For each representation, rank features by the configuration's importance method and select the top $b_k$, after rounding $b^*_k$ as in ``From base representations to a hybrid''.
    \item Join the selected features to form the hybrid representation.
\end{enumerate}

% TODO: REVIEW MODIFICATION
Each fixed-share method follows the steps described in ``From base representations to a hybrid''. It rescales each standalone score to $\tilde r_k$, sets $b_k$ from $w_k = g(\tilde r_k)$ and takes the top $b_k$ features of representation $k$. Allocation runs only over the representations retained by the representation count. The standalone scores cost one fit per representation, fold and internal evaluator, and every budget, allocation method and importance method reuses them.

% TODO: review, needs manual review
This method cannot adapt allocation based on feature complementarity, and it relies on the assumption that many features in a well-performing representation are valuable.

% TODO: REVIEW MODIFICATION
Fixed shares rely on the assumption that many features in a well-performing representation are valuable. This is not always the case, since a strong representation may owe its score to a few features, so that a large share fills with weak ones. Fixed shares also cannot adapt the allocation to feature complementarity, and two representations carrying the same information both receive a share. Of the five methods, only the global pool ranks features from different representations against each other.

% TODO: explain why the greedy approach was tried, but it's too computationally intensive and actually quite ineffective
% TODO: explain framework experiments and how they helped here

% CODE INPUT
The four fixed-share methods differ only in the weight $w_k$ computed from the rescaled score $\tilde r_k$ (Table~\ref{tab:alloc_methods}), and their settings are fixed for every dataset rather than searched. Linear allocation sets $w_k = \tilde r_k$ and squared allocation sets $w_k = \tilde r_k^{\,2}$, so under both, a representation receives 0 features when its rescaled score is 0. Softmax allocation sets $w_k = \exp\big((\tilde r_k - \max_l \tilde r_l)/T\big)$, with temperature $T = 1$. Rank-decay allocation sets $w_k = e^{-0.3\rho_k}$, where $\rho_k$ is the representation's rank by score, 0 for the best. Under these two methods, no representation receives 0 features because of its score. Since every rescaled score lies between 0 and 1, softmax allocation cannot give one representation most of the budget, and with 10 representations the best takes at most 23.2\% of it. Under rank decay only the order counts, so with 10 representations the best receives 27.3\% and the tenth 1.8\%, whatever the gaps between their scores. Table~\ref{tab:alloc_example} works through one example with a budget of 100 features and three representations with rescaled scores 0.6, 0.5 and 0.2.

% CODE INPUT
\begin{table}[htbp]
\centering
\caption[The five allocation methods searched.]{The five allocation methods searched. For each method, the columns give the weight $w_k$ computed from representation $k$'s rescaled standalone score $\tilde r_k$, the setting fixed for every dataset, how features are chosen and when a representation receives no features. $B$ is the budget in features, $b_k$ is representation $k$'s share of the budget, $\rho_k$ is its rank by score and $T$ is the softmax temperature. Under every method, a representation also receives no features when the representation count leaves it out or its standalone fit fails.}
\label{tab:alloc_methods}
\small
\begin{adjustbox}{max width=\textwidth}
\begin{tabular}{@{}P{0.13\linewidth}P{0.24\linewidth}P{0.14\linewidth}P{0.24\linewidth}P{0.19\linewidth}@{}}
\toprule
Method & Weight $w_k$ & Fixed setting & How features are chosen & When a representation receives 0 features \\
\midrule
Linear & $\tilde r_k$ & none & Top $b_k$ of its own ranking & $\tilde r_k = 0$ \\
Squared & $\tilde r_k^{\,2}$ & power 2 & Top $b_k$ of its own ranking & $\tilde r_k = 0$ \\
Softmax & $\exp\big((\tilde r_k - \max_l \tilde r_l)/T\big)$ & $T = 1$ & Top $b_k$ of its own ranking & never from its score \\
Rank decay & $e^{-0.3\rho_k}$, $\rho_k$ = rank by score, 0 for the best & decay 0.3 & Top $b_k$ of its own ranking & never from its score \\
Global pool & $\tilde r_k^{\,2}$, sizing its offer to a pool of $2B$ candidates & power 2; pool $2B$ & Top $B$ of one joint ranking over all offered candidates & $\tilde r_k = 0$, or none of its candidates in the joint top $B$ \\
\bottomrule
\end{tabular}
\end{adjustbox}
\end{table}

% CODE INPUT
\begin{table}[htbp]
\centering
\caption[Worked example of the four fixed-share allocation methods.]{Worked example of the four fixed-share allocation methods, with a budget of 100 features split among three representations. Each column is one representation, headed by its rescaled standalone score $\tilde r_k$, and each entry is the number of features it receives. Shares are rounded by largest remainder, as described in ``From base representations to a hybrid''.}
\label{tab:alloc_example}
\begin{adjustbox}{max width=\textwidth}
\begin{tabular}{@{}lccc@{}}
\toprule
Method & $\tilde r_1 = 0.6$ & $\tilde r_2 = 0.5$ & $\tilde r_3 = 0.2$ \\
\midrule
Linear & 46 & 39 & 15 \\
Squared & 55 & 39 & 6 \\
Softmax & 39 & 35 & 26 \\
Rank decay & 44 & 32 & 24 \\
\bottomrule
\end{tabular}
\end{adjustbox}
\end{table}

% CODE INPUT
The global pool sizes each representation's offer with squared weights $\tilde r_k^{\,2}$, as shares of a pool of $2B$ candidate features. Each offer is capped at the representation's width in features, and only features scoring above 0 are offered. The offered features are joined into one matrix, on which the same importance method is fitted once, on the internal training part. The hybrid keeps the top $B$ features of this joint ranking, again only those scoring above 0. Each representation's final count is therefore an outcome of the joint ranking rather than set in advance, and there is no minimum of one feature per representation. A representation receives 0 features when its score is 0 or when none of its candidates reach the joint top $B$. Under the global pool, a feature's rank depends on which features the other representations offered, whereas under fixed shares it does not. The cost is one extra importance fit per configuration, on the pooled candidates.

% CODE INPUT
Under all five methods, a representation also receives no features when the representation count leaves it out or when its standalone fit fails.

% CODE INPUT
We first tried greedy forward selection, which picks features for how well they work together, something fixed shares cannot do \citep{Kohavi1997, Guyon2003}. In each round, it added the next block of features from whichever representation most improved the internal validation score. It stopped when the budget was spent or when the score had not improved for a set number of rounds. With blocks of $s$ features and $K$ representations, it needs up to $(B/s) \times K$ fits per configuration and fold. With blocks of 10 features, this amounts to 500 fits at 500 features and 10 representations, against 10 for the score-based methods. In preliminary runs, it reached a Spearman correlation of 0.3600 on Jain antibody FcRn retention, against 0.4526 for linear allocation, and placed 16th of 17 methods. On breast cancer PAM50 subtype, it reached an accuracy of 0.8551 against 0.8938 for linear allocation, again placing 16th of 17, and it was removed from the search.

% CODE INPUT
The five methods were chosen in a preliminary comparison run before the main experiments. It compared 8 candidate allocation methods, each at 4 representation counts, on 9 regression datasets. It used 3 downstream models (LightGBM, random forest and CatBoost) and budgets of 500 and 1,000 features, giving 54 combinations of dataset, downstream model and budget. A fixed-share method was the single best method in 38 of the 54 combinations and the global pool in 16. All five were kept because no single method was best in most combinations. This comparison predates the current score rescaling, and it sized the weights with an inner 3-fold group split instead of the 80/20 internal split used now.

\subsubsection{Search configuration and internal evaluation}
% [notes moved to thesis_notes_moved.md]

% TODO: review, needs manual review
Features that improve predictions under RF may not necessarily help a kNN model, and vice versa. This is because variable importance depends fundamentally on the model, as Dong and Rudin showed \citep{Dong2020}, and as we found in our feature importance scoring analysis. To address this, we use an internal evaluator which scores each base representation on its own with a model whose inductive bias best fits the downstream model. We use four evaluators, namely RF, DNN, LightGBM and the matched evaluator, a special case that matches the internal evaluator to the final downstream model. We fit each on the internal training part of the training fold and score it on the internal validation part.

% TODO: REVIEW MODIFICATION
Features that improve predictions under a random forest may not necessarily help an SVM, and vice versa, because variable importance depends fundamentally on the model, as Dong and Rudin showed \citep{Dong2020}. To address this, we use an internal evaluator, a model fitted inside the search on one representation at a time and separate from the downstream model scored at the end. We search four evaluators, namely random forest, DNN, LightGBM and the matched evaluator, a special case that matches the internal evaluator to the downstream model. No evaluator is chosen in advance to suit the downstream model. Instead, all four are searched and the best configuration wins. Features are standardized on the internal training part of the training fold, where we fit each evaluator before scoring it on the internal validation part in the dataset's own metric.

% TODO: review, needs manual review
This internal evaluator is used to score each representation, and those scores set how the budget is split and which representations are kept within the larger hybrid search. Thus, it is another setting we sweep while searching for the best set of features.

% TODO: REVIEW MODIFICATION
The internal evaluator scores each representation, and those scores set how the budget is split and which representations are kept within the larger hybrid search. It does not rank features inside a representation, which is the task of the importance method. Judging a set of features by the score of a model fitted on it is the wrapper idea of \citet{Kohavi1997}. The internal evaluator is thus another setting swept in the search for the best set of features.

% TODO: need additional files to specify architectures

% TODO: does match always work? custom downstream model?

% TODO: these numbers are outdated
% TODO: review, needs manual review
The full search tests all combinations of budgets, allocations, importance methods, representation counts, downstream models, and internal evaluators under 5-fold cross-validation: 8,160 configurations per dataset when using 5 downstream models. We approach this problem combinatorially as it allows us to measure which factors contribute most to performance variance (model choice, budget, allocation, or their interactions) rather than assuming any single axis dominates.

% TODO: REVIEW MODIFICATION
The full search tests all combinations of budgets, allocation methods, importance methods, representation counts, downstream models and internal evaluators under 5-fold cross-validation. This gives 8,160 configurations per dataset where all five downstream models run, fewer where a dataset leaves a downstream model out, and 10,080 on humanness, where a Gaussian process is added. A search at the fast tier has 816 configurations. The breakdown by downstream model and the number of model fits are given in Additional file~\ref{app:add:search_size}. We search every combination rather than assume that one setting dominates. This lets us measure which settings, or which interactions between them, account for most of the spread in scores.

% TODO: mention representation knobs

% TODO: review, needs manual review
The complete search result is cached to disk. The search runs once; everything after that is a lookup.

% TODO: REVIEW MODIFICATION
The complete search result is cached to disk. Importance scores are computed once per fold and importance method, and each representation's standalone score once per fold and internal evaluator. Every configuration reuses them, so the cost grows more slowly than the number of configurations. Because each configuration's score is stored by its settings, a fast search can later be extended to a full one. The chosen configuration for each pair of downstream model and budget is saved as a blueprint. The search thus runs once, and every later analysis is a lookup.

When I began to implement this idea, I figured that there would be a fairly straightforward way to select the most informative features. However, I began to ask some questions. How do you measure feature importance? If you're starting with a handful of input representations, how many features do you take from each? How do you score these potential feature sets?  Unfortunately, my one-size-fits-all hypothesis rang true on multiple fronts. As I found out, different prediction tasks prefer different answers to all of those questions. So to figure out how to even go about selecting the best features, I realized I had to sweep a search space.

% TODO: explain framework experiments
We ran a series of experiments to determine the best variables to include in the search space.
When making a selection, we took into account how many potential hybrid representations would be gained by including that variable, as well as the compute time

% TODO: REVIEW MODIFICATION
Before the full runs, we ran a series of preliminary design experiments to determine the values each setting of the search may take. They used 1 to 9 datasets and 1 to 5 folds, and each dataset was scored in its own metric, with scores never averaged across datasets. When making a selection, we considered whether a value gave the best score for some dataset, downstream model and budget, as well as the compute time in hours. A value that won somewhere but cost many hours stayed in the full search and was left out of the fast tier. Each setting's values were then either kept, because each won for some dataset, model and budget, or fixed at one value. A setting was fixed when the spread between its values was smaller than the spread of one value across folds. As an example of a kept value, nine downstream models were fitted on each of 10 base representations over 5 folds. The datasets were Buchwald--Hartwig yield, Caco-2 efflux, Koenig expression, LogD and breast cancer PAM50 subtype. The DNN scored best on 4 of the 5 datasets and third on Caco-2 (0.2936 against 0.2851 RMSE, log$_{10}$ efflux ratio), so it was added. As an example of a fixed rule, we tested five ways of rescaling each representation's score before the budget is split. On hERG $K_i$, with linear allocation at 300 features, the five rescalings gave RMSE values from 0.5870 to 0.5924 (p$K_i$). A single rescaling, however, ranged from 0.5475 to 0.6294 RMSE across the 5 folds, so rescaling is one fixed rule per metric. Every design experiment is listed in Additional file~\ref{app:add:design_experiments}.

% CODE INPUT
We search four values of the internal evaluator, namely the matched evaluator, random forest, LightGBM and DNN, and the fast tier searches the same four. The matched evaluator is an evaluator of the same model family as the downstream model. For a random forest, LightGBM or DNN downstream model, the matched evaluator is already in the list, so these three downstream models are searched with 3 evaluators. For an XGBoost or SVM downstream model, the matched evaluator adds an XGBoost or SVM evaluator, so these two are searched with 4. On humanness, where a Gaussian process is also a downstream model, the matched evaluator adds a Gaussian process evaluator, so it is searched with 4 as well. The same family does not imply the same settings, since only the DNN evaluator shares the downstream DNN's architecture and training settings. Even for the DNN, one difference remains. The evaluator draws the 20\% of molecules it uses for early stopping at random, whereas the downstream DNN splits its 20\% by the same groups as the outer split. The tree evaluators are smaller than the downstream trees, with 50 trees against 100. Their depth is limited to 10 for the random forest and to 4 for LightGBM and XGBoost, which keeps scoring every representation inexpensive. The SVM and Gaussian process evaluators, whose cost grows with the square to the cube of the number of molecules, are fitted on at most 500 training molecules drawn with a fixed seed. The random forest evaluator tries all features at each split on regression datasets and the square root of the number of features on classification datasets. The downstream random forest tries the square root on both. If the DNN evaluator fails to fit on a fold, a random forest scores that fold and the run records a warning. Every evaluator setting is given in Additional file~\ref{app:add:model_settings}, alongside the settings of the same model used as a downstream model.

% CODE INPUT
A user can also supply their own downstream model under any name, provided it can be fitted and can predict. The matched evaluator, however, knows only a fixed list of model families, so a model outside that list receives no matched evaluator, and no warning is given. That model is then scored only by the other evaluators, by default random forest, LightGBM and DNN. If the matched evaluator was the only one requested, a random forest scores the representations. A hand-written blueprint that specifies the matched evaluator and is rebuilt with a model the search does not know is scored by LightGBM instead. This is a limitation of the method, and a user with their own model should name an evaluator explicitly.

% CODE INPUT
Variable importance depends on the model \citep{Dong2020}, so the representation that a random forest scores highest need not be the one an SVM uses best. Matching the evaluator to the downstream model is the obvious default. However, the evaluator only has to rank representations, and another model may rank them better for the downstream model. Because the matched evaluator is one of the four searched values, the search already covers this default, and a mismatched evaluator winning is informative rather than a negative result. We first counted wins in the search itself, in a preliminary run on 5 datasets (nanobody thermostability, hERG $K_i$, humanness, LogD and breast cancer PAM50 subtype). These runs used 10 representations, budgets of 25 to 300 features and 5 folds, and their scores are the search's own fold scores rather than held-out scores. The matched evaluator gave the best score for its downstream model in 9 of 25 pairs of dataset and downstream model. Counted per budget, it was best in 43 of 125 combinations of dataset, downstream model and budget. With 3 or 4 evaluators per downstream model and no effect of matching, about 38 of 125 would be expected. On nanobody thermostability with a DNN downstream model, the matched evaluator reached a Spearman correlation of 0.6696 against 0.6506 for the LightGBM evaluator, a gap of 0.0190. Its 5 fold scores ranged from 0.5095 to 0.7920. On LogD with an SVM downstream model, the matched evaluator scored 0.4322 RMSE (log units) against 0.4217 for the LightGBM evaluator, where lower is better. This is a gap of 0.0105, the matched evaluator lost at all 5 budgets, and its 5 fold scores ranged from 0.4118 to 0.4634. We then scored the evaluators on held-out folds in a second preliminary run, with 3 outer folds, a budget of 250 features, 5 representations and LightGBM gain only. The downstream models were random forest and SVM. On nanobody thermostability with a random forest downstream model, the matched evaluator reached a Spearman correlation of 0.6098, the LightGBM evaluator 0.6220 and the SVM evaluator 0.6172. The matched evaluator's fold scores ranged from 0.5659 to 0.6734. With an SVM downstream model, the matched evaluator reached 0.6189, the Gaussian process evaluator 0.6265 and the LightGBM evaluator 0.6317, with fold scores for the matched evaluator from 0.5530 to 0.6703. On hERG $K_i$ with a random forest downstream model, the matched evaluator scored 0.6046 RMSE (p$K_i$), the SVM evaluator 0.6059 and the LightGBM evaluator 0.6075. Its fold scores ranged from 0.5957 to 0.6178, and it beat the SVM evaluator on 2 of 3 folds. With an SVM downstream model, the matched evaluator scored 0.5723 RMSE (p$K_i$), the LightGBM evaluator 0.5724 and the Gaussian process evaluator 0.5765, with fold scores for the matched evaluator from 0.5605 to 0.5862. Every gap above, both held out and within the search, is smaller than the range of the matched evaluator's own fold scores. \textbf{[RESULT NEEDED: per dataset and downstream model, how often the winning configuration used the matched evaluator, and its gap to the best other evaluator in the dataset's metric]}


\subsection{Augmentation}
% [notes moved to thesis_notes_moved.md]

% TODO: review, also I think this is out of date
% TODO: review, needs manual review
Some features are too low-dimensional or too specialized to take part in the budget allocation stage, for instance MACCS fingerprints (167 features) or log-likelihood ratios (3 features). Excluding them from the allocation would lose potentially useful signal, whereas including them would distort the budget distribution because they cannot contribute many features. We therefore handle these features as augmentation. After the base hybrid is constructed, we pool the augmentation features and rank them by how strongly each correlates with the out-of-fold errors of a small random forest fitted on the base features. We keep the top 0, 5, 10, 20 or 40 of them, whichever scores best with the downstream model on the internal validation set.

% TODO: REVIEW MODIFICATION
Some features are too low-dimensional or too specialized to take part in the budget allocation stage, for instance MACCS fingerprints (167 features) or log-likelihood ratios (3 features). Others, such as mutation features that record differences from a parent antibody, only carry meaning alongside a full representation. Excluding such features from the allocation would lose potentially useful signal, whereas including them would distort the budget distribution because they cannot contribute many features. We therefore handle them as augmentation. We call the finished hybrid $\mathbf{H}$ the base, and we call the feature sets from which the extra features are drawn the donors. After the base is constructed, we pool the donor features and rank them by how strongly each correlates with the out-of-fold errors of a small random forest fitted on the base features. We keep the top 0, 5, 10, 20 or 40 of them, whichever scores best with the downstream model on the internal validation set. In the uncapped mode, the kept features are appended outside the budget. In the capped mode, they replace the weakest base features, so that the hybrid stays within the budget.

% TODO: I need a better list to describe augmentations
Augmentation features range from smaller fingerprints like MACCS or Coloumb matrices, sets of structural information like 3D pharmacophores, and features describing only the differences between the sames like mutation features for antibodies.

% TODO: REVIEW MODIFICATION
Donors are set per dataset rather than per domain and are listed in the Augmentation column of Table~\ref{tab:dataset_reps}. Augmentation features range from smaller fingerprints such as MACCS to sets of structural information such as 3D pharmacophores and Coulomb matrix eigenvalues. They also include features that describe only the differences between variants, such as mutation features for antibodies. The Coulomb matrix donor is not a fingerprint but the 50 largest eigenvalues of a matrix built from atomic numbers and 3D distances between atoms. MACCS is usually a donor, serving as one on four of the six small-molecule datasets, and it is a base representation on MDR1--MDCK efflux. QM9 is the exception, since MACCS is not offered there and the topological torsion fingerprint (2,048 features) is a donor, the only wide one. Two datasets, humanness and trastuzumab--HER2 binding, have no donors, and nanobody thermostability has one, mutation features. Each donor is small, at most 167 features apart from the torsions, or describes only one aspect such as shape, charge, the mutated positions or one tumor property. Each is therefore a plausible supplement but a weak representation on its own. Evidence from the literature suggests that 3D donors are worth offering. \citet{Bahia2023} modeled six protein--ligand series (461 structures) with 2D, 3D and combined 2D and 3D descriptors, using $k$-nearest neighbors, random forest and lasso. Combining 2D and 3D descriptors gave many more significant models, which the authors attribute to the two descriptor types coding for complementary properties. Their 3D descriptors came from bioactive conformations in the Protein Data Bank, whereas our 3D donors use one generated conformer per molecule.
% TODO: add augmentation subsection

% CODE INPUT
Donor features are ranked by what the base gets wrong. Inside each outer training fold, an 80/20 internal split is made with each scaffold or sequence group kept on one side, and the outer test fold is never used. A small random forest (50 trees, depth at most 10) is fitted on the base features, and its predictions for the internal-training molecules come from 3-fold group-aware cross-validation. Its out-of-fold residuals are $r_i = y_i - \hat{y}^{\mathrm{oof}}_i$, where for binary tasks $\hat{y}$ is the predicted probability of the positive class. We use out-of-fold residuals because a forest's in-sample residuals shrink toward zero and mostly measure fit noise. All donor features are pooled, and each is scored by its absolute Pearson correlation with the residuals, $c_j = |\rho_{\mathrm{P}}(\mathbf{x}_j, \mathbf{r})|$. The top $M = 50$ features proceed to the next step. Ranking against the labels $\mathbf{y}$ would favor features that repeat what the base already captures, whereas ranking against the residuals $\mathbf{r}$ favors features that explain its errors. The same idea underlies gradient boosting, which fits each new stage to the current residual \citep{friedman2001greedy}, and least angle regression, which adds the variable most correlated with the residual \citep{efron2004lars}. The residual model is always this inexpensive forest and never the downstream model, so the ranking costs the same for every downstream model.

% CODE INPUT
Block sizes of $k \in \{0, 5, 10, 20, 40\}$ features are tried, capped at the number of candidates, where block $k$ consists of the top $k$ ranked features. Each block is added to the base and scored with the downstream model on the internal-validation molecules. The chosen size $k^*$ has the best score and must beat $k=0$. The option of no extras, $k=0$, is scored first and wins ties, so keeping no extras is the default outcome rather than a failure. Five whole-block choices replace about 50 one-by-one decisions to accept or reject a feature. Choosing among many options on a small validation set overfits the choice \citep{Cawley2010}, and five options leave less room for this. The chosen donors and feature positions are stored with the rest of the search settings and are re-extracted for the outer fold like the rest of the hybrid. \textbf{[RESULT NEEDED: per dataset and downstream model, how often $k^*=0$ is chosen, and per-fold scores for no augmentation, uncapped and capped]}

% CODE INPUT
There are two ways to fit the extras against the budget $B$. In uncapped mode, extras are appended outside the budget, so the final width is $|\mathbf{H}| + k^*$, at most $B + 40$ features. In capped mode, extras must displace base features. The base features are ranked once by the importance of a random forest (50 trees, depth at most 10), and the $k^*$ weakest are dropped, so the final width stays at most $B$. In capped mode, any unspent budget is filled first, and at least one base feature is always kept. Only in capped mode does an extra feature have to beat a base feature to earn its place. We run both modes because uncapped augmentation and no augmentation differ in width at the same nominal $B$. Their comparison therefore mixes better features with more features, and capped mode removes this difference in width.

% CODE INPUT
Augmentation is a separate condition rather than a searched setting. For each dataset with donors, the whole search runs three times, with no augmentation, uncapped augmentation and capped augmentation. Every other setting (downstream model, budget, representation count, allocation method, importance method and internal evaluator) is searched within each condition. Fourteen of the 16 datasets run all three conditions, while humanness and trastuzumab--HER2 binding have no donors and so run without augmentation only. Without augmentation, the width is at most $B$ and nothing is offered. Uncapped augmentation allows at most $B + 40$ features, and an extra feature must beat keeping no extras. Capped augmentation keeps the width at most $B$, and an extra feature must beat both keeping no extras and the weakest base feature it replaces. The three conditions are also run under the nearest-neighbor search objective of Chapter~\ref{chap:similarity}. Two datasets never run that objective. Breast cancer PAM50 subtype has 4 classes, whereas the objective is scored by AUROC, which needs 2. On trastuzumab--HER2 binding, a held-out block can contain only one group. Augmentation also has a second use outside the main search, in which a user keeps one strong representation at full width as the base and adds a few extras. In that case $B$ is the width of that representation, so only uncapped mode makes sense. We test this use in a separate experiment that adds a few features to one strong representation, with decoy features planted to check that the rule admits none of them. \textbf{[RESULT NEEDED: how many decoy features are admitted in the decoy check]} This use matters for adoption, because a chemist can keep the representation they already trust and learn which few outside features help, and how many. Augmentation also asks directly whether another representation holds signal the base lacks (Chapter~\ref{chap:flaws}). On Koenig expression, the mutation features and the ESM-2 log-likelihood ratios reach the model only through augmentation. \textbf{[RESULT NEEDED: which donors are kept most often, per domain]}

% CODE INPUT
The residual ranking has four limitations. First, $|\rho_{\mathrm{P}}|$ captures only linear relations with the residual, one feature at a time. Second, it ignores redundancy among the extras, so two near-identical donor features can both rank high. Third, donor features are pooled with no limit per donor, so a 2,048-feature donor gets more chances in the ranking than a 49-feature one. Fourth, the margin over $k=0$ is zero, so a near-tie on one internal split can admit extras.

% CODE INPUT
Table~\ref{tab:aug_conditions} summarizes the three augmentation conditions.

% CODE INPUT
\begin{table}[htbp]
\centering
\small
\caption[The three augmentation conditions.]{The three augmentation conditions. For each condition, the columns give the final width of the hybrid, what an extra feature must beat to be kept and the datasets that run the condition. $B$ is the budget and $|\mathbf{H}|$ the width of the base hybrid, both in features, and $k^*$ is the chosen number of extra features. An extra feature is one drawn from a donor, a small or specialized feature set offered to the hybrid.}
\label{tab:aug_conditions}
\begin{adjustbox}{max width=\textwidth}
\begin{tabular}{@{}P{2.4cm}P{3.0cm}P{4.6cm}P{3.6cm}@{}}
\toprule
Condition & Final width (features) & What an extra feature must beat & Datasets that run it \\
\midrule
No augmentation & at most $B$ & nothing is offered & all 16 \\
Uncapped & $|\mathbf{H}| + k^*$, at most $B + 40$ & keeping no extras & the 14 with donors \\
Capped & at most $B$ & keeping no extras, and the weakest base feature it replaces & the 14 with donors \\
\bottomrule
\end{tabular}
\end{adjustbox}
\end{table}


\subsection{Fast tier and adaptive pruning}
% [notes moved to thesis_notes_moved.md]

% CODE INPUT
The search can be made cheaper in two ways, which differ in when they decide. The fast tier uses shorter lists of values for some settings, which are fixed before the search and identical on every fast-tier dataset. Adaptive pruning, in contrast, decides during the search from the dataset's own fold scores. The fast tier knowingly accepts a possibly worse answer in exchange for a much smaller grid. Pruning instead aims to find the same best configuration while no longer scoring configurations that have already clearly lost. Without either, the full search does not finish on the largest datasets.

% CODE INPUT
Six of the 16 datasets use the fast tier, namely QM9, LogD, Koenig expression, trastuzumab--HER2 binding, H3K4me3 and TATA promoter. They are marked in Table~\ref{tab:dataset_reps}, and all other datasets use the full search. A dataset uses the fast tier when the full search would take longer than about 40 hours, the time that four consecutive 12-hour jobs can cover. It also uses the fast tier when it has more than 10,000 molecules or sequences. The estimated run times of the full search were about 103 hours for QM9, 64 for Koenig expression, 62 for LogD and 39 for TATA promoter, and H3K4me3 has 18,244 sequences. The values the fast tier drops are still searched on every other dataset. The fast tier is thus a compromise on compute and makes no claim that those values are useless. A fast run can later be extended to a full one without redoing work, because each configuration's score is stored by its settings rather than by the tier. A full search then scores only the configurations the fast one skipped.

% CODE INPUT
The fast lists were chosen by one rule in a preliminary comparison. A value stays in the fast lists if it is inexpensive and improves the best score, whereas an expensive value that wins often stays in the full search only. Importance methods and downstream models were measured on five datasets, with 5 folds each, against the unpruned full search. The five were breast cancer PAM50 subtype, nanobody thermostability, humanness, hERG $K_i$ and LogD. DeepLIFT took about half of the fast run's hours on all five datasets (49.9\% to 51.6\%), and dropping it changed only LogD's best score, from 0.4905 to 0.4930 RMSE (log units). Adding the DNN as a downstream model moved LogD from 0.4930 to 0.4273 RMSE, against a full-search best of 0.4170. Adding the SVM moved hERG $K_i$ from 0.5776 to 0.5599 RMSE (p$K_i$), while adding XGBoost changed no dataset's best score. All five downstream models are now searched at both tiers. The global pool won 60 of the 125 combinations of dataset, downstream model and budget on those five datasets. It is expensive, however, so it is in the full search only. The three fast allocation methods are linear, squared and rank-decay allocation. \textbf{[RESULT NEEDED: the share of run time the fast tier and adaptive pruning saved on the full runs]}

\subsubsection{Search tier settings}
% [notes moved to thesis_notes_moved.md]

% TODO: review, needs manual review
Table~\ref{tab:search_tiers} lists, for each search setting, the values tried at the full tier and at the fast tier.

% TODO: review, needs manual review
\begin{table}[htbp]
\centering
\caption[Values searched for each setting at the two search tiers.]{Values searched for each setting at the full and fast search tiers. Both tiers use five folds, and configurations removed by adaptive pruning are scored on the first two folds only. The matched evaluator is an evaluator of the same model family as the downstream model. SVM = support vector machine; DNN = deep neural network.}
\label{tab:search_tiers}
\begin{adjustbox}{max width=\textwidth}
\begin{tabular}{@{}P{0.22\linewidth}P{0.40\linewidth}P{0.30\linewidth}@{}}
\toprule
Setting & Full tier & Fast tier \\
\midrule
Downstream models & Random forest, XGBoost, LightGBM, SVM$^{a}$, DNN; Gaussian process$^{b}$ & Random forest, XGBoost, LightGBM, SVM$^{a}$, DNN \\
Internal evaluators & Matched evaluator, random forest, LightGBM, DNN & Same as full \\
Feature budgets (features) & 10, 25, 50, 100, 250, 500, 1{,}000, 2{,}000 & Same as full \\
Allocation methods & Linear, squared, softmax, rank decay, global pool & Linear, squared, rank decay \\
Importance methods & LightGBM gain, DeepLIFT, TreeSHAP & LightGBM gain \\
Representation counts & 2, 3, 5, all representations & 3, all representations \\
\bottomrule
\end{tabular}
\end{adjustbox}

\smallskip
{\footnotesize $^{a}$SVM is not fitted above 10{,}000 molecules (5{,}000 when augmentation features are added outside the budget). $^{b}$Added only on Humanness. Of the 16 datasets, it is the only one that is under 1{,}000 molecules, has its size recorded and is not excluded for run time.}
\end{table}

\subsubsection{Adaptive pruning}
% [notes moved to thesis_notes_moved.md]

% TODO: review, needs manual review
In adaptive pruning, every configuration is first scored on the first two folds. At a budget of 100 features, within each group of configurations that share a downstream model, internal evaluator and representation count, pairs of allocation and importance method with no configuration in the better half are not searched at larger budgets. Each setting is then voted on within groups that fix every other setting. An allocation method is pruned if it is in the best 2 in fewer than 20\% of groups, and a representation count if it is best in fewer than 10\%. An importance method is pruned if it is in the best 2 in fewer than 20\%, and an internal evaluator if it is in the best 3 in fewer than 10\%. Survivors are scored on the remaining folds.

% TODO: REVIEW MODIFICATION
In adaptive pruning, every configuration is first scored on the first two folds, and the rules rank configurations by their mean score over those two folds, in the dataset's own metric. Budgets are scored from smallest to largest. Once every budget up to 100 features has been scored, a checkpoint runs within each group of configurations that share a downstream model, internal evaluator and representation count. Within each such group, pairs of allocation and importance method with no configuration in the better half are not searched at larger budgets, although at least one pair per group is kept. Each setting is then voted on within groups that fix every other setting, including the budget. The best few values in each group receive a point, and a value is removed if it receives a point in fewer than a cut-off share of groups. An allocation method is pruned if it is among the best 2 in fewer than 20\% of groups, and a representation count if it is the best in fewer than 10\%. An importance method is pruned if it is among the best 2 in fewer than 20\% of groups, and an internal evaluator if it is among the best 3 in fewer than 10\%. A vote runs only if the setting has more values than receive a point, and it can never remove every value of a setting. A configuration without a score is ranked last in its group. The surviving configurations are scored on the remaining folds, and the best configuration is picked by its mean over all five folds, exactly as without pruning.

% CODE INPUT
Some configurations are so far behind after two folds that the other three folds would only confirm a loss, and adaptive pruning stops scoring them. Pruning is judged on one criterion, namely that it must never remove a configuration that would have been the best, and the time saved is secondary. Pruning is off unless a search requests it. As submitted, the search stored once per dataset uses pruning on every dataset, at both the fast and the full tier. On a fast-tier dataset only one importance method is searched, so the importance vote does not run. Hyperparameter search uses a related idea, dropping configurations early on partial evaluations, as in successive halving \citep{Jamieson2016} and Hyperband \citep{Li2018Hyperband}. Here the partial evaluation is two of the five folds.

% CODE INPUT
Table~\ref{tab:pruning_rules} lists the five pruning rules.

% CODE INPUT
\begin{table}[htbp]
\centering
\small
\caption[The adaptive pruning rules.]{The adaptive pruning rules. For each rule, the columns give how groups of configurations are formed, which values are kept or receive a point, and the cut-off below which a value is removed. The last column gives what the rule removed on the five calibration datasets in a preliminary run. Every rule ranks configurations by their mean score over folds 1 and 2, in the dataset's own metric.}
\label{tab:pruning_rules}
\begin{adjustbox}{max width=\textwidth}
\begin{tabular}{@{}P{2.2cm}P{3.4cm}P{2.8cm}P{1.6cm}P{3.6cm}@{}}
\toprule
Rule & How groups are formed & Kept, or receives a point & Cut-off & Removed on the calibration datasets \\
\midrule
Checkpoint at 100 features & configurations at budgets up to 100 features with the same downstream model, internal evaluator and representation count & (allocation, importance) pairs with a configuration in the better half of the group & better half of the configurations & configurations within groups only; never a whole allocation or importance method \\
\addlinespace
Representation count & every other setting fixed, budget included & best 1 & 10\% of groups & the count of 5 on LogD; the count of 2 on breast cancer PAM50 subtype \\
\addlinespace
Allocation method & every other setting fixed, budget included & best 2 & 20\% of groups & softmax on LogD \\
\addlinespace
Importance method & every other setting fixed, budget included & best 2 & 20\% of groups & DeepLIFT on nanobody thermostability and hERG $K_i$; TreeSHAP on LogD \\
\addlinespace
Internal evaluator & every other setting fixed, budget included & best 3 & 10\% of groups & the SVM evaluator on LogD \\
\bottomrule
\end{tabular}
\end{adjustbox}
\end{table}

% CODE INPUT
Because the votes are grouped by every other setting, values compete only with each other. A value that is best for a random forest at 25 features is therefore not removed for losing to a DNN at 500 features. A single fold is not enough for the votes. In the representation-count vote on breast cancer PAM50 subtype after fold 1 alone, 726 of 1,020 groups had every value at an identical score. Ties are then broken by list order, so a vote after one fold removes values by accident. After one fold, no cut-off of any vote was safe on the calibration datasets, whereas after two folds each vote had at least one safe cut-off.

% CODE INPUT
The cut-offs were calibrated in a preliminary study on five datasets, LogD, hERG $K_i$, breast cancer PAM50 subtype, nanobody thermostability and humanness. Each dataset's grid was first scored without pruning on all five folds. Each rule was then replayed on the stored fold scores, so that no rule needed new fits. A cut-off counts as safe if the best five-fold score is unchanged for each of the 25 combinations of downstream model and budget, on every dataset. Where several safe cut-offs made the same cuts, we chose the one whose neighboring cut-off was also safe, and the saving in run time was not a criterion. The chosen rules removed the count of 5 representations on LogD and the count of 2 on PAM50 subtype. They also removed softmax allocation on LogD, DeepLIFT on nanobody thermostability and hERG $K_i$, TreeSHAP on LogD and the SVM evaluator on LogD. The checkpoint removed configurations within groups only and never a whole allocation or importance method. TreeSHAP accounts for 71\% of LogD's fitting time and was never the best on LogD, so most of the saving comes from there. Estimated run times for eight budgets, without and with pruning, were 79 and 44 hours on LogD, 42 and 35 on hERG $K_i$, and 42 and 30 on PAM50 subtype. On nanobody thermostability they were 32 and 24 hours, and on humanness 15 and 13. The calibration applied the votes one after another, whereas the search applies all four at once to the same scores from folds 1 and 2. The cut-offs were therefore checked in a different order from the one the search uses. Five datasets form a small calibration set, which is a limitation of the rules.

% CODE INPUT
Pruning changes what the reported score means. With pruning on, folds 1 and 2 both select the survivors and score them, so the five-fold mean of the winner is slightly optimistic \citep{Cawley2010}. Single representations, by contrast, are never pruned. A fair comparison of the hybrid against a single representation therefore leaves both voting folds out of both scores, or re-scores without them.

